% ─────────────────────────────────────────────────────────────────────────────
%  Resumen de Probabilidad — Modelos Estocásticos
%  Autor: P. Hurtado · 2026
% ─────────────────────────────────────────────────────────────────────────────
\documentclass[10pt,a4paper]{article}

\usepackage[T1]{fontenc}
\usepackage[utf8]{inputenc}
\usepackage{lmodern}
\usepackage[a4paper, top=1.8cm, bottom=2.0cm, left=1.5cm, right=1.5cm]{geometry}
\usepackage[spanish,es-noshorthands]{babel}
\usepackage{amsmath,amssymb,mathtools,amsthm}
\usepackage{xcolor}
\usepackage[most]{tcolorbox}
\usepackage{tikz}
\usepackage{pgfplots}
\pgfplotsset{compat=1.18}
\usepackage{microtype}
\usepackage{needspace}
\usepackage{parskip}
\usepackage{fancyhdr}
\usepackage{hyperref}
\usepackage{booktabs}
\usepackage{array}
\usepackage{caption}

\usetikzlibrary{arrows.meta, positioning, calc, shapes.geometric, fit, backgrounds, decorations.pathreplacing}

\setlength{\parskip}{3pt}
\setlength{\parindent}{0pt}
\setlength{\headheight}{19pt}

% ─── Numeración de ecuaciones por sección ────────────────────────────────────
\numberwithin{equation}{section}

% ─── Colores ─────────────────────────────────────────────────────────────────
\definecolor{sectionbg}{rgb}{0.16,0.26,0.52}
\definecolor{subsecbg}{rgb}{0.38,0.50,0.75}
\definecolor{codeblue}{rgb}{0.0,0.28,0.68}
\definecolor{notebg}{rgb}{1.0,0.97,0.87}
\definecolor{noteborder}{rgb}{0.80,0.65,0.25}
\definecolor{defbg}{rgb}{0.93,0.98,0.93}
\definecolor{defborder}{rgb}{0.25,0.60,0.30}
\definecolor{thmbg}{rgb}{0.92,0.95,1.0}
\definecolor{thmborder}{rgb}{0.25,0.45,0.80}
\definecolor{exbg}{rgb}{0.96,0.96,0.96}
\definecolor{exborder}{rgb}{0.55,0.55,0.55}
\definecolor{warnbg}{rgb}{1.0,0.94,0.93}
\definecolor{warnborder}{rgb}{0.80,0.25,0.20}
\definecolor{proofbg}{rgb}{0.98,0.97,1.0}
\definecolor{proofborder}{rgb}{0.55,0.40,0.75}

% ─── Hyperref ────────────────────────────────────────────────────────────────
\hypersetup{
  colorlinks=true,
  linkcolor=codeblue,
  citecolor=codeblue,
  urlcolor=codeblue,
}

% ─── Cabeceras ───────────────────────────────────────────────────────────────
\pagestyle{fancy}
\fancyhf{}
\fancyhead[L]{\footnotesize\sffamily\color{sectionbg}Resumen de Probabilidad --- Modelos Estoc\'asticos}
\fancyhead[R]{\footnotesize\sffamily\color{sectionbg}2026}
\fancyfoot[C]{\footnotesize\sffamily\thepage}
\renewcommand{\headrulewidth}{0.4pt}
\renewcommand{\headrule}{\hbox to\headwidth{\color{subsecbg}\leaders\hrule height \headrulewidth\hfill}}

% ─── tcolorbox styles ────────────────────────────────────────────────────────
\tcbset{
  defbox/.style={
    enhanced,
    colback=defbg, colframe=defborder,
    boxrule=0.6pt, arc=3pt,
    left=7pt, right=7pt, top=5pt, bottom=5pt,
    fonttitle=\small\bfseries\sffamily,
    coltitle=white, colbacktitle=defborder,
    titlerule=0pt,
    colupper=black,
  },
  thmbox/.style={
    enhanced,
    colback=thmbg, colframe=thmborder,
    boxrule=0.6pt, arc=3pt,
    left=7pt, right=7pt, top=5pt, bottom=5pt,
    fonttitle=\small\bfseries\sffamily,
    coltitle=white, colbacktitle=thmborder,
    titlerule=0pt,
  },
  exbox/.style={
    enhanced,
    colback=exbg, colframe=exborder,
    boxrule=0.4pt, arc=3pt,
    left=7pt, right=7pt, top=5pt, bottom=5pt,
    fonttitle=\small\bfseries\sffamily,
    coltitle=white, colbacktitle=exborder,
    titlerule=0pt,
  },
  notebox/.style={
    colback=notebg, colframe=noteborder,
    boxrule=0.4pt, arc=3pt,
    left=5pt, right=5pt, top=4pt, bottom=4pt,
    fontupper=\footnotesize\itshape,
  },
  warnbox/.style={
    enhanced,
    colback=warnbg, colframe=warnborder,
    boxrule=0.6pt, arc=3pt,
    left=7pt, right=7pt, top=5pt, bottom=5pt,
    fonttitle=\small\bfseries\sffamily,
    coltitle=white, colbacktitle=warnborder,
    titlerule=0pt,
    title={$\triangle$\ Atenci\'on},
    fontupper=\small,
  },
  proofbox/.style={
    enhanced,
    colback=proofbg, colframe=proofborder!60,
    boxrule=0.6pt, arc=3pt,
    left=7pt, right=7pt, top=5pt, bottom=5pt,
    fonttitle=\small\itshape\sffamily,
    coltitle=white, colbacktitle=proofborder!80,
    titlerule=0pt,
    title={Demostraci\'on},
    fontupper=\small,
  },
}

% ─── Helpers ─────────────────────────────────────────────────────────────────
\newcommand{\secheader}[2]{%
  \refstepcounter{section}%
  \addcontentsline{toc}{section}{\protect\numberline{\thesection}#2}%
  \vspace{10pt}\needspace{6cm}\noindent%
  \begin{tcolorbox}[
    colback=sectionbg, colupper=white, colframe=sectionbg,
    sharp corners, left=8pt, right=8pt, top=5pt, bottom=5pt,
    fontupper=\large\bfseries\sffamily,
  ]
  Secci\'on #1:\enspace #2
  \end{tcolorbox}\vspace{2pt}%
}

\newcommand{\subsecheader}[1]{%
  \needspace{4cm}\vspace{6pt}\noindent%
  \addcontentsline{toc}{subsection}{#1}%
  \begin{tcolorbox}[
    colback=subsecbg!18, colframe=subsecbg!45,
    sharp corners, boxrule=0.4pt,
    left=5pt, right=5pt, top=2pt, bottom=2pt,
    fontupper=\small\bfseries\sffamily,
  ]
  #1
  \end{tcolorbox}\vspace{1pt}%
}

% ─── Macros matemáticos ──────────────────────────────────────────────────────
\newcommand{\Prob}[1]{\mathbb{P}\!\left(#1\right)}
\newcommand{\E}[1]{\mathbb{E}\!\left[#1\right]}
\newcommand{\Var}[1]{\mathrm{Var}\!\left(#1\right)}
\newcommand{\Cov}[2]{\mathrm{Cov}\!\left(#1,\,#2\right)}
\newcommand{\sigalg}{\mathcal{F}}
\newcommand{\ind}{\mathbf{1}}
\newcommand{\Borel}{\mathcal{B}}
\newcommand{\R}{\mathbb{R}}
\newcommand{\Nset}{\mathbb{N}}
\DeclareMathOperator{\supp}{supp}

% ─── Contadores ──────────────────────────────────────────────────────────────
\newcounter{defcount}[section]
\renewcommand{\thedefcount}{\thesection.\arabic{defcount}}
\newcounter{thmcount}[section]
\renewcommand{\thethmcount}{\thesection.\arabic{thmcount}}
\newcounter{excount}[section]
\renewcommand{\theexcount}{\thesection.\arabic{excount}}

% ─── Entornos ────────────────────────────────────────────────────────────────
\newenvironment{definicion}[1]{%
  \refstepcounter{defcount}%
  \begin{tcolorbox}[defbox, title={Definici\'on \thedefcount\ \textperiodcentered\ #1}]%
}{\end{tcolorbox}}

\newenvironment{teorema}[1]{%
  \refstepcounter{thmcount}%
  \begin{tcolorbox}[thmbox, title={Teorema \thethmcount\ \textperiodcentered\ #1}]%
}{\end{tcolorbox}}

\newenvironment{proposicion}[1]{%
  \refstepcounter{thmcount}%
  \begin{tcolorbox}[thmbox, title={Proposici\'on \thethmcount\ \textperiodcentered\ #1}]%
}{\end{tcolorbox}}

\newenvironment{ejemplo}[1]{%
  \refstepcounter{excount}%
  \begin{tcolorbox}[exbox, title={Ejemplo \theexcount\ \textperiodcentered\ #1}]%
}{\end{tcolorbox}}

\newenvironment{nota}{\begin{tcolorbox}[notebox]}{\end{tcolorbox}}

\newenvironment{cuidado}{\begin{tcolorbox}[warnbox]}{\end{tcolorbox}}

\newenvironment{demostracion}{\begin{tcolorbox}[proofbox]}{\hfill$\square$\end{tcolorbox}}

% ─────────────────────────────────────────────────────────────────────────────
\begin{document}
% ─────────────────────────────────────────────────────────────────────────────

%  T I T U L O
\begin{tcolorbox}[
    colback=sectionbg, colupper=white, colframe=sectionbg,
    sharp corners, left=12pt, right=12pt, top=10pt, bottom=10pt,
  ]
  {\Large\bfseries\sffamily Resumen de Probabilidad}\\[4pt]
  {\normalsize\sffamily Modelos Estoc\'asticos}\\[6pt]
  {\small\sffamily\color{white!80!sectionbg}
    Espacio muestral \textperiodcentered\ Probabilidad \textperiodcentered\ Variables aleatorias \textperiodcentered\
    Distribuciones \textperiodcentered\ Vectores aleatorios \textperiodcentered\ Proceso de Poisson}\\[6pt]
  {\footnotesize\sffamily\color{white!85!sectionbg}\itshape
    Por P. Hurtado \textperiodcentered\ 2026}
\end{tcolorbox}

\vspace{4pt}
\begin{nota}
  Estas notas siguen el enfoque medida-te\'orico de Kolmogorov. Se asume familiaridad con
  an\'alisis real (l\'imites, series, integrales de Lebesgue a nivel introductorio).
  Notaci\'on: $\Omega$ espacio muestral, $\sigalg$ $\sigma$-\'algebra, $P$ medida de probabilidad,
  $X$ variable aleatoria.
\end{nota}

\tableofcontents

\vspace{8pt}
\begin{tcolorbox}[
    colback=white, colframe=sectionbg,
    boxrule=0.6pt, arc=3pt,
    left=8pt, right=8pt, top=6pt, bottom=6pt,
    fonttitle=\small\bfseries\sffamily,
    title={Leyenda de colores},
  ]
  \small
  Cada bloque de contenido se identifica visualmente por color, seg\'un su naturaleza:
  \vspace{4pt}

  \renewcommand{\arraystretch}{1.35}
  \begin{tabular}{@{}p{0.7cm} l p{10.5cm}@{}}
    \tikz\fill[defbg, draw=defborder, line width=0.6pt] (0,0) rectangle (0.55,0.35);
                                          & \textbf{Definici\'on}      & Concepto formal con nombre y notaci\'on. Marco verde claro.                            \\
    \tikz\fill[thmbg, draw=thmborder, line width=0.6pt] (0,0) rectangle (0.55,0.35);
                                          & \textbf{Teorema/Proposici\'on} & Resultado demostrable a partir de las definiciones. Marco azul.                    \\
    \tikz\fill[proofbg, draw=proofborder!50, line width=0.6pt] (0,0) rectangle (0.55,0.35);
                                          & \textbf{Demostraci\'on}    & Argumento que justifica un teorema/proposici\'on. Marco morado.                        \\
    \tikz\fill[exbg, draw=exborder, line width=0.6pt] (0,0) rectangle (0.55,0.35);
                                          & \textbf{Ejemplo}           & Aplicaci\'on concreta o caso ilustrativo. Marco gris.                                  \\
    \tikz\fill[notebg, draw=noteborder, line width=0.6pt] (0,0) rectangle (0.55,0.35);
                                          & \textbf{Nota}              & Observaci\'on, intuici\'on o comentario complementario. Marco \'ambar.                 \\
    \tikz\fill[warnbg, draw=warnborder, line width=0.6pt] (0,0) rectangle (0.55,0.35);
                                          & \textbf{Atenci\'on}        & Advertencia sobre un error com\'un o sutileza. Marco rojo.                             \\
    \tikz\fill[sectionbg] (0,0) rectangle (0.55,0.35);
                                          & \textbf{Secci\'on}         & Encabezado principal del m\'odulo. Fondo azul oscuro.                                  \\
    \tikz\fill[subsecbg!18, draw=subsecbg!45, line width=0.6pt] (0,0) rectangle (0.55,0.35);
                                          & \textbf{Subsecci\'on}      & Subdivisi\'on tem\'atica dentro de una secci\'on. Fondo azul claro.                    \\
  \end{tabular}
\end{tcolorbox}

\newpage

% =============================================================================
\secheader{1}{Espacio de Probabilidad $(\Omega,\,\sigalg,\,P)$}

\subsecheader{1.1 \textperiodcentered\ Experimento, espacio muestral y eventos}

\begin{definicion}{Experimento aleatorio}
  Un \textbf{experimento aleatorio} es todo procedimiento u observaci\'on cuyo resultado
  no puede predecirse con certeza antes de realizarlo, pero que cumple:
  \begin{enumerate}
    \item El conjunto de resultados posibles es conocido \emph{a priori}.
    \item Cada realizaci\'on (\emph{ensayo} o \emph{prueba}) produce exactamente un resultado.
    \item Es \emph{repetible} bajo condiciones esencialmente id\'enticas, exhibiendo regularidad
          estad\'istica en frecuencias relativas a largo plazo.
  \end{enumerate}
  Se contrasta con un \textbf{experimento determinista}, en el que las mismas condiciones
  iniciales producen siempre el mismo resultado.
\end{definicion}

\begin{ejemplo}{Experimentos aleatorios t\'ipicos}
  Lanzar una moneda, tirar un dado, medir el tiempo de espera de un cliente en una cola,
  contar el n\'umero de llamadas a un call center en una hora, observar el precio de cierre
  diario de una acci\'on.
\end{ejemplo}

\begin{definicion}{Espacio muestral}
  El \textbf{espacio muestral} $\Omega$ asociado a un experimento aleatorio es el conjunto
  de todos los resultados posibles de dicho experimento. Cada elemento $\omega \in \Omega$
  se llama \emph{resultado} o \emph{punto muestral}.\\[4pt]
  Un \textbf{evento} $A$ es cualquier subconjunto $A \subseteq \Omega$.
\end{definicion}

\begin{ejemplo}{Lanzamiento de dos dados}
  $\Omega = \{(i,j) : i,j \in \{1,2,3,4,5,6\}\}$, $|\Omega|=36$.\\
  Evento ``suma igual a 7'': $A = \{(1,6),(2,5),(3,4),(4,3),(5,2),(6,1)\}$, $|A|=6$.\\
  Por la regla cl\'asica: $P(A) = 6/36 = 1/6$.
\end{ejemplo}

\begin{definicion}{Operaciones entre eventos}
  Sean $A, B \subseteq \Omega$ eventos. Se definen:
  \begin{itemize}
    \item \textbf{Uni\'on}: $A \cup B = \{\omega \in \Omega : \omega \in A \;\text{o}\; \omega \in B\}$
          \hfill (ocurre $A$ o $B$, o ambos).
    \item \textbf{Intersecci\'on}: $A \cap B = \{\omega \in \Omega : \omega \in A \;\text{y}\; \omega \in B\}$
          \hfill (ocurren ambos).
    \item \textbf{Complemento}: $A^c = \{\omega \in \Omega : \omega \notin A\}$
          \hfill (no ocurre $A$).
    \item \textbf{Diferencia}: $A \setminus B = A \cap B^c$
          \hfill (ocurre $A$ pero no $B$).
  \end{itemize}
  $A$ y $B$ son \textbf{disjuntos} (mutuamente excluyentes) si $A \cap B = \emptyset$.\\[3pt]
  \textbf{Leyes de De Morgan}: $(A\cup B)^c = A^c \cap B^c$, \quad $(A\cap B)^c = A^c \cup B^c$.
\end{definicion}

\vspace{4pt}
\begin{center}
  \begin{tikzpicture}[scale=0.9]
    \draw[thick, rounded corners=4pt, fill=thmbg!60] (-4,-2.5) rectangle (4,2.5);
    \node[above left, font=\small\bfseries, color=sectionbg] at (3.8,2.4) {$\Omega$};
    \draw[fill=defbg, draw=defborder, thick] (-1.2,0) ellipse (1.6cm and 1.5cm);
    \node[font=\small\bfseries, color=defborder] at (-2.0,0.3) {$A$};
    \draw[fill=notebg, draw=noteborder, thick] (1.2,0) ellipse (1.6cm and 1.5cm);
    \node[font=\small\bfseries, color=noteborder] at (2.0,0.3) {$B$};
    \node[font=\tiny\bfseries] at (0,0) {$A\cap B$};
    \node[font=\small, color=gray] at (-3.2, 2.0) {$A^c\cap B^c$};
    \node[font=\small, color=gray] at (-2.5,-0.0) {$A\setminus B$};
    \node[font=\small, color=gray] at (2.5, 0.0) {$B\setminus A$};
  \end{tikzpicture}
\end{center}
\begin{center}
  {\footnotesize\itshape Diagrama de Venn: relaciones entre eventos $A$, $B$ dentro de $\Omega$.}
\end{center}

\begin{definicion}{Conjunto potencia}
  El \textbf{conjunto potencia} de $\Omega$, denotado $2^\Omega$ o $\mathcal{P}(\Omega)$, es la
  colecci\'on de todos los subconjuntos de $\Omega$:
  \[
    2^\Omega = \mathcal{P}(\Omega) = \{A : A \subseteq \Omega\}.
  \]
  En particular $\emptyset \in 2^\Omega$ y $\Omega \in 2^\Omega$.\\[3pt]
  Si $\Omega$ es finito con $|\Omega| = n$, entonces $|2^\Omega| = 2^n$ (de ah\'i la notaci\'on).
  El conjunto potencia es el universo m\'as grande posible de eventos sobre $\Omega$ y constituye
  siempre una $\sigma$-\'algebra (la m\'as fina).
\end{definicion}

\subsecheader{1.2 \textperiodcentered\ Sigma-\'algebra}

Una \textbf{$\sigma$-\'algebra} (o \emph{tribu}) es la estructura matem\'atica que selecciona
\emph{cu\'ales} subconjuntos de $\Omega$ admiten asignaci\'on de probabilidad. Intuitivamente,
es el \emph{cat\'alogo de eventos observables}: la colecci\'on de aquellos subconjuntos sobre
los que la pregunta ``\textit{\textquestiondown cu\'al es la probabilidad de que ocurra?}''
tiene respuesta bien definida.

\begin{nota}
  \textbf{\textquestiondown Por qu\'e no usar simplemente $2^\Omega$?} Cuando $\Omega$ es no numerable
  (p.\,ej.\ $\Omega=\R$), no existe una medida de probabilidad consistente sobre \emph{todos} los
  subconjuntos: existen conjuntos no medibles (conjuntos de Vitali, paradoja de Banach--Tarski).
  La $\sigma$-\'algebra restringe el dominio a eventos para los cuales $P$ s\'i puede definirse
  de forma coherente con los axiomas de Kolmogorov.\\[3pt]
  Para $\Omega$ finito o numerable s\'i se usa $\sigalg = 2^\Omega$ sin problemas.
\end{nota}

Formalmente, una $\sigma$-\'algebra es una familia de subconjuntos cerrada bajo las operaciones
b\'asicas (complemento y uniones \emph{numerables}), que adem\'as contiene al espacio total. El
prefijo ``$\sigma$-'' enfatiza la cerradura bajo operaciones \emph{numerables} (no solo finitas):
esto es lo que permite manejar l\'imites y series, esenciales en probabilidad.

\begin{definicion}{$\sigma$-\'algebra}
  Una familia $\sigalg \subseteq 2^\Omega$ es una \textbf{$\sigma$-\'algebra} si:
  \begin{enumerate}
    \item $\Omega \in \sigalg$
    \item $A \in \sigalg \Rightarrow A^c \in \sigalg$ \hfill (cerrada bajo complemento)
    \item $\{A_n\}_{n\geq 1} \subseteq \sigalg \Rightarrow \bigcup_{n=1}^{\infty} A_n \in \sigalg$
          \hfill (cerrada bajo uniones numerables)
  \end{enumerate}
  Los elementos de $\sigalg$ se llaman \textbf{eventos medibles}.
\end{definicion}

\begin{nota}
  Se derivan: $\emptyset \in \sigalg$; cerrada bajo intersecciones numerables; cerrada bajo
  diferencias. La $\sigma$-\'algebra m\'as peque\~na es $\{\emptyset,\Omega\}$; la m\'as grande
  es $2^\Omega$ (conjunto potencia).
\end{nota}

\subsecheader{1.3 \textperiodcentered\ Medida de probabilidad}

\begin{definicion}{Medida de probabilidad (Kolmogorov, 1933)}
  Una funci\'on $P : \sigalg \to [0,1]$ es una \textbf{medida de probabilidad} si:
  \begin{align}
     & \textbf{(K1)}\quad P(A) \geq 0 \quad \forall A \in \sigalg \tag{no negatividad}              \\
     & \textbf{(K2)}\quad P(\Omega) = 1 \tag{normalizaci\'on}                                       \\
     & \textbf{(K3)}\quad P\!\left(\bigsqcup_{n=1}^{\infty} A_n\right) = \sum_{n=1}^{\infty} P(A_n)
    \quad \text{para } A_i \cap A_j = \emptyset \; \forall i\neq j \tag{$\sigma$-aditividad}
  \end{align}
  La terna $(\Omega, \sigalg, P)$ es el \textbf{espacio de probabilidad}.
\end{definicion}

\begin{nota}
  \textbf{$\sigma$-aditividad vs.\ partici\'on.} La $\sigma$-aditividad (K3) aplica a
  \emph{cualquier} colecci\'on disjunta $\{A_n\}$; su uni\'on no tiene por qu\'e ser $\Omega$.
  Una \textbf{partici\'on} es un caso particular: una familia disjunta que adem\'as
  \emph{cubre} todo $\Omega$, es decir, $\bigsqcup_n A_n = \Omega$. En ese caso,
  $\sum_n P(A_n) = P(\Omega) = 1$.
\end{nota}

\vspace{2pt}
\begin{center}
  \begin{tikzpicture}[scale=0.85]
    \begin{scope}[shift={(-4.5,0)}]
      \draw[thick, rounded corners=4pt, fill=thmbg!40] (-2.5,-1.8) rectangle (2.5,1.8);
      \node[font=\small\bfseries, color=sectionbg] at (-2.2,1.55) {$\Omega$};
      \draw[fill=defbg, draw=defborder, thick] (-1.4,0.4) ellipse (0.7cm and 0.55cm);
      \node[font=\small\bfseries, color=defborder] at (-1.4,0.4) {$A_1$};
      \draw[fill=defbg, draw=defborder, thick] (0.2,-0.5) ellipse (0.7cm and 0.55cm);
      \node[font=\small\bfseries, color=defborder] at (0.2,-0.5) {$A_2$};
      \draw[fill=defbg, draw=defborder, thick] (1.5,0.6) ellipse (0.7cm and 0.55cm);
      \node[font=\small\bfseries, color=defborder] at (1.5,0.6) {$A_3$};
      \node[font=\footnotesize\itshape, color=gray] at (-1.7,-1.4) {resto: no cubierto};
      \node[font=\small\bfseries, color=sectionbg] at (0,-2.4) {$\sigma$-aditividad};
      \node[font=\footnotesize, color=black] at (0,-2.95)
      {$P\!\left(\bigsqcup_n A_n\right)=\sum_n P(A_n)\leq 1$};
    \end{scope}
    \begin{scope}[shift={(4.5,0)}]
      \draw[thick, rounded corners=4pt, fill=thmbg!40] (-2.5,-1.8) rectangle (2.5,1.8);
      \node[font=\small\bfseries, color=sectionbg] at (-2.2,1.55) {$\Omega$};
      \fill[defbg, draw=defborder, thick] (-2.5,-1.8) -- (-2.5,1.8) -- (-0.8,1.8) -- (-1.2,-1.8) -- cycle;
      \fill[notebg, draw=noteborder, thick] (-0.8,1.8) -- (-1.2,-1.8) -- (0.5,-1.8) -- (0.3,1.8) -- cycle;
      \fill[exbg, draw=exborder, thick] (0.3,1.8) -- (0.5,-1.8) -- (1.6,-1.8) -- (1.4,1.8) -- cycle;
      \fill[proofbg, draw=proofborder!60, thick] (1.4,1.8) -- (1.6,-1.8) -- (2.5,-1.8) -- (2.5,1.8) -- cycle;
      \node[font=\small\bfseries, color=defborder] at (-1.75,0) {$B_1$};
      \node[font=\small\bfseries, color=noteborder] at (-0.4,0) {$B_2$};
      \node[font=\small\bfseries, color=exborder] at (0.95,0) {$B_3$};
      \node[font=\small\bfseries, color=proofborder] at (2.0,0) {$B_4$};
      \node[font=\small\bfseries, color=sectionbg] at (0,-2.4) {Partici\'on};
      \node[font=\footnotesize, color=black] at (0,-2.95)
      {$\bigsqcup_n B_n=\Omega \;\Rightarrow\; \sum_n P(B_n)=1$};
    \end{scope}
  \end{tikzpicture}
\end{center}
\begin{center}
  {\footnotesize\itshape Izquierda: eventos disjuntos cuya uni\'on no llena $\Omega$ (basta para
    aplicar K3). Derecha: partici\'on, donde los disjuntos cubren todo $\Omega$ y sus
    probabilidades suman 1.}
\end{center}

\begin{proposicion}{Propiedades b\'asicas}
  Para todo $A, B \in \sigalg$:
  \begin{enumerate}
    \item $P(\emptyset) = 0$
    \item $P(A^c) = 1 - P(A)$
    \item $A \subseteq B \Rightarrow P(A) \leq P(B)$
    \item $P(A \cup B) = P(A) + P(B) - P(A \cap B)$ \quad (inclusi\'on y exclusi\'on)
    \item $P(A) \in [0,1]$
  \end{enumerate}
\end{proposicion}

\begin{nota}
  \textbf{T\'ecnica clave.} Las cinco propiedades son consecuencias \emph{directas} de los
  axiomas K1--K3. El truco recurrente es \emph{descomponer un evento en piezas disjuntas}
  y aplicar la $\sigma$-aditividad (K3) sobre esa descomposici\'on.
\end{nota}

\begin{demostracion}
  \textbf{(1)} $P(\emptyset)=0$. \quad
  Escribir $\Omega = \Omega \sqcup \emptyset \sqcup \emptyset \sqcup \cdots$.
  Por K3: $P(\Omega) = P(\Omega) + \sum_{n\geq 1} P(\emptyset)$.
  Restando $P(\Omega)$ (finito por K2) y usando K1 ($P(\emptyset)\geq 0$): $P(\emptyset) = 0$.
  \\[5pt]
  \textbf{(2)} $P(A^c)=1-P(A)$. \quad
  $\Omega = A \sqcup A^c$ (disjuntos), luego por K3: $1 = P(A)+P(A^c)$. Despejar.
  \\[5pt]
  \textbf{(3)} $A\subseteq B \Rightarrow P(A)\leq P(B)$. \quad
  Si $A\subseteq B$, entonces $B = A \sqcup (B\setminus A)$ (disjuntos). Por K3:
  $P(B) = P(A) + P(B\setminus A) \geq P(A)$, pues $P(B\setminus A)\geq 0$ por K1.
  \\[5pt]
  \textbf{(4)} $P(A\cup B) = P(A)+P(B)-P(A\cap B)$. \quad
  Descomponer ambos lados en piezas disjuntas:
  \begin{align*}
    A\cup B & = A \sqcup (B\setminus A)         &  & \Rightarrow\; P(A\cup B) = P(A)+P(B\setminus A), \\
    B       & = (A\cap B) \sqcup (B\setminus A) &  & \Rightarrow\; P(B\setminus A) = P(B)-P(A\cap B).
  \end{align*}
  Sustituyendo: $P(A\cup B) = P(A) + P(B) - P(A\cap B)$.
  \\[5pt]
  \textbf{(5)} $P(A)\in[0,1]$. \quad
  $P(A)\geq 0$ por K1. Como $A\subseteq\Omega$, por (3) y K2: $P(A)\leq P(\Omega)=1$.
\end{demostracion}

\begin{teorema}{Principio de inclusi\'on y exclusi\'on}
  Sea $\{A_1, \ldots, A_n\} \subseteq \sigalg$ una colecci\'on finita de eventos. Entonces:
  \[
    P\!\left(\bigcup_{i=1}^{n} A_i\right)
    \;=\; \sum_{k=1}^{n} (-1)^{k+1}
    \sum_{\substack{S\subseteq\{1,\ldots,n\}\\|S|=k}} P\!\left(\bigcap_{i\in S} A_i\right).
  \]
  Expandiendo los primeros t\'erminos:
  \[
    P\!\left(\bigcup_i A_i\right) = \sum_i P(A_i) \;-\; \sum_{i<j} P(A_i\cap A_j)
    \;+\; \sum_{i<j<k} P(A_i\cap A_j\cap A_k) \;-\; \cdots
    \;+\; (-1)^{n+1} P(A_1\cap\cdots\cap A_n).
  \]
\end{teorema}

\begin{nota}
  \textbf{Idea intuitiva.} Sumar $P(A_1)+\cdots+P(A_n)$ \emph{sobreestima} $P(\bigcup A_i)$
  porque las intersecciones se cuentan varias veces. Restando los pares ($\sum P(A_i\cap A_j)$)
  se compensa, pero ahora las triples quedan mal contadas; se vuelven a sumar, etc.
  El signo alterna y los t\'erminos se cancelan exactamente.
\end{nota}

\begin{demostracion}
  Inducci\'on sobre $n$.
  \\[5pt]
  \textbf{Caso base $n=2$.} Es la propiedad (4) ya demostrada:
  $P(A_1\cup A_2) = P(A_1)+P(A_2)-P(A_1\cap A_2)$.
  \\[5pt]
  \textbf{Paso inductivo $(n-1)\Rightarrow n$.}
  Sea $B = \bigcup_{i=1}^{n-1} A_i$ y supongamos la f\'ormula v\'alida para $n-1$ eventos.
  Por el caso base aplicado a $B$ y $A_n$:
  \begin{equation}
    P(B \cup A_n) \;=\; P(B) \,+\, P(A_n) \,-\, P(B \cap A_n). \tag{$\star$}
  \end{equation}
  Por hip\'otesis de inducci\'on (HI):
  \[
    P(B) = \sum_{k=1}^{n-1} (-1)^{k+1}
    \sum_{\substack{S\subseteq\{1,\ldots,n-1\}\\|S|=k}} P\!\left(\bigcap_{i\in S} A_i\right).
  \]
  Por distributividad: $B\cap A_n = \bigcup_{i=1}^{n-1}(A_i \cap A_n)$. Aplicando la HI
  a los $n-1$ eventos $\{A_i \cap A_n\}_{i=1}^{n-1}$:
  \[
    P(B\cap A_n) = \sum_{k=1}^{n-1} (-1)^{k+1}
    \sum_{\substack{S\subseteq\{1,\ldots,n-1\}\\|S|=k}}
    P\!\left(A_n \cap \bigcap_{i\in S} A_i\right).
  \]
  Sustituyendo en $(\star)$ y reagrupando t\'erminos por el tama\~no $k$ del subconjunto
  de \'indices, se obtiene la f\'ormula para $n$.
\end{demostracion}

\begin{nota}
  \textbf{Desigualdades de Bonferroni.} Las sumas parciales del lado derecho alternan como
  cotas: tomando solo el primer t\'ermino se obtiene la cota superior (subaditividad)
  $P(\bigcup A_i)\leq \sum P(A_i)$; con dos t\'erminos se obtiene la cota inferior
  $P(\bigcup A_i)\geq \sum P(A_i) - \sum_{i<j} P(A_i\cap A_j)$, y as\'i sucesivamente.
\end{nota}

\subsecheader{1.4 \textperiodcentered\ Recap de combinatoria}

Cuando $\Omega$ es finito con resultados \emph{equiprobables}, la \textbf{regla cl\'asica} entrega:
\[
  P(A) = \frac{|A|}{|\Omega|}, \qquad A \subseteq \Omega.
\]
El c\'alculo de probabilidades se reduce entonces a \emph{contar}: cu\'antos resultados favorables
y cu\'antos posibles.

\begin{definicion}{Principio multiplicativo}
  Si un procedimiento se realiza en $k$ etapas y la etapa $i$-\'esima admite $n_i$ resultados
  posibles (independientes de las anteriores), el n\'umero total de resultados es:
  \[
    n_1 \cdot n_2 \cdots n_k = \prod_{i=1}^{k} n_i.
  \]
\end{definicion}

\begin{ejemplo}{Placas patentes}
  Una patente de 4 letras seguidas de 2 d\'igitos admite $26^4 \cdot 10^2 = 45{,}697{,}600$
  combinaciones distintas (con repetici\'on permitida).
\end{ejemplo}

\begin{definicion}{Permutaciones y variaciones}
  De un conjunto de $n$ elementos distintos:
  \begin{itemize}
    \item \textbf{Permutaci\'on de los $n$ elementos} (orden total):
          $P_n = n! = n\cdot(n-1)\cdots 2\cdot 1$.
    \item \textbf{Variaciones sin repetici\'on} (secuencias ordenadas de $k\leq n$):
          \[
            V(n,k) = \frac{n!}{(n-k)!} = n(n-1)(n-2)\cdots(n-k+1).
          \]
    \item \textbf{Variaciones con repetici\'on} (secuencias de longitud $k$, repetici\'on permitida):
          $V_R(n,k) = n^k$.
  \end{itemize}
\end{definicion}

\begin{definicion}{Combinaciones}
  Subconjuntos de tama\~no $k$ tomados de un conjunto de $n$ elementos (\emph{el orden no importa}):
  \[
    \binom{n}{k} = \frac{n!}{k!\,(n-k)!}, \qquad 0 \leq k \leq n.
  \]
  \textbf{Propiedades}:
  \begin{itemize}
    \item Simetr\'ia: $\binom{n}{k} = \binom{n}{n-k}$.
    \item Bordes: $\binom{n}{0} = \binom{n}{n} = 1$.
    \item Regla de Pascal: $\binom{n}{k} = \binom{n-1}{k-1} + \binom{n-1}{k}$.
    \item Teorema del binomio: $(x+y)^n = \sum_{k=0}^{n}\binom{n}{k} x^k y^{n-k}$.
  \end{itemize}
  \textbf{Combinaciones con repetici\'on} (multiconjuntos de tama\~no $k$ con $n$ tipos disponibles):
  \[
    CR(n,k) = \binom{n+k-1}{k} \qquad (\textit{stars and bars}).
  \]
\end{definicion}

\begin{nota}
  \textbf{\textquestiondown Qu\'e f\'ormula usar?} Antes de calcular, identifica dos preguntas:
  \begin{itemize}
    \item \textbf{\textquestiondown Importa el orden?} S\'i $\to$ variaci\'on/permutaci\'on. \quad No $\to$ combinaci\'on.
    \item \textbf{\textquestiondown Se permite repetir?} S\'i $\to$ con repetici\'on. \quad No $\to$ sin repetici\'on.
  \end{itemize}
\end{nota}

\begin{center}
  \renewcommand{\arraystretch}{1.45}
  \begin{tabular}{@{}l c c l@{}}
    \toprule
    \textbf{Tipo}                  & \textbf{Orden} & \textbf{Repetici\'on} & \textbf{F\'ormula}      \\
    \midrule
    Permutaci\'on de $n$           & s\'i           & no                    & $n!$                    \\
    Variaci\'on $V(n,k)$           & s\'i           & no                    & $\dfrac{n!}{(n-k)!}$    \\
    Variaci\'on con repetici\'on   & s\'i           & s\'i                  & $n^k$                   \\
    Combinaci\'on $\binom{n}{k}$   & no             & no                    & $\dfrac{n!}{k!\,(n-k)!}$ \\
    Combinaci\'on con repetici\'on & no             & s\'i                  & $\binom{n+k-1}{k}$      \\
    \bottomrule
  \end{tabular}
\end{center}

\begin{ejemplo}{Aplicaci\'on: regla cl\'asica + combinatoria}
  \textbf{Planteamiento.} De una baraja est\'andar de 52 cartas se reparte una mano
  de 5 cartas al azar. \textquestiondown Cu\'al es la probabilidad de que la mano
  contenga exactamente los 4 ases?
  \\[5pt]
  \textbf{(1) Definir el experimento y el espacio muestral.}
  Cada ``mano'' es un subconjunto de $5$ cartas escogidas de $52$. Como la mano se
  identifica solo por \emph{qu\'e} cartas contiene (el orden de reparto no influye en
  las reglas del juego), trabajamos con \textbf{combinaciones}, no variaciones:
  \[
    |\Omega| = \binom{52}{5} = \frac{52!}{5!\,47!} = 2\,598\,960
    \quad\text{(manos posibles, sin orden, sin repetici\'on).}
  \]
  Asumimos que el reparto es uniforme: cada una de esas $\binom{52}{5}$ manos tiene la
  misma probabilidad. Aplica entonces la regla cl\'asica $P(A) = |A|/|\Omega|$.
  \\[5pt]
  \textbf{(2) Definir el evento.} Sea $A = \{$la mano contiene los 4 ases$\}$.
  Una mano en $A$ se construye en dos pasos independientes:
  \begin{itemize}
    \item Elegir los $4$ ases: hay $\binom{4}{4}=1$ forma (los cuatro ases son obligatorios).
    \item Elegir la $5^a$ carta entre las $48$ no-ases restantes: $\binom{48}{1}=48$ formas.
  \end{itemize}
  Por el principio multiplicativo:
  \[
    |A| = \binom{4}{4}\binom{48}{1} = 1 \cdot 48 = 48.
  \]
  \\[-5pt]
  \textbf{(3) Calcular la probabilidad.}
  \[
    P(A) = \frac{|A|}{|\Omega|} = \frac{48}{2\,598\,960} \approx 1.85\times 10^{-5}.
  \]
  \\[-5pt]
  \textbf{(4) Interpretaci\'on.} Aproximadamente $1$ de cada $54{,}145$ manos contiene
  los cuatro ases.
\end{ejemplo}

% =============================================================================
\secheader{2}{Probabilidad Condicional e Independencia}

\subsecheader{2.1 \textperiodcentered\ Probabilidad condicional}

\begin{definicion}{Probabilidad condicional}
  Para $A, B \in \sigalg$ con $P(B) > 0$:
  \[
    P(A \mid B) \;=\; \frac{P(A \cap B)}{P(B)}.
  \]
  La funci\'on $P(\,\cdot \mid B) : \sigalg \to [0,1]$ es en s\'i misma una medida de probabilidad
  sobre $(\Omega, \sigalg)$.
\end{definicion}

\begin{nota}
  Intuitivamente: al conocer que ocurri\'o $B$, el espacio muestral efectivo se reduce a $B$
  y se re-normaliza dividiendo por $P(B)$.
\end{nota}

\begin{teorema}{Ley de probabilidad total}
  Sea $\{B_1, \ldots, B_n\}$ una partici\'on de $\Omega$ con $P(B_i) > 0$. Entonces:
  \[
    P(A) = \sum_{i=1}^{n} P(A \mid B_i)\, P(B_i) \qquad \forall\, A\in\sigalg.
  \]
\end{teorema}

\begin{demostracion}
  Como $\{B_1, \ldots, B_n\}$ es una partici\'on de $\Omega$:
  \[
    A \;=\; A \cap \Omega \;=\; A \cap \!\left(\bigsqcup_{i=1}^{n} B_i\right)
    \;=\; \bigsqcup_{i=1}^{n} (A \cap B_i),
  \]
  donde la \'ultima igualdad usa la distributividad de $\cap$ sobre $\sqcup$, y los $A\cap B_i$
  son disjuntos porque los $B_i$ lo son. Por $\sigma$-aditividad (K3):
  \[
    P(A) \;=\; \sum_{i=1}^{n} P(A \cap B_i).
  \]
  Por la definici\'on de probabilidad condicional, $P(A\cap B_i) = P(A\mid B_i)\,P(B_i)$
  (v\'alida pues $P(B_i)>0$). Sustituyendo:
  \[
    P(A) \;=\; \sum_{i=1}^{n} P(A \mid B_i)\, P(B_i).
  \]
\end{demostracion}

\begin{teorema}{Teorema de Bayes}
  Sean $A, B \in \sigalg$ con $P(A)>0$ y $P(B)>0$. Entonces:
  \[
    P(B \mid A) \;=\; \frac{P(A \mid B)\, P(B)}{P(A)}.
  \]
  Si $P(A)$ no se conoce directamente, se obtiene por probabilidad total con la partici\'on
  $\{B, B^c\}$:
  \[
    P(B \mid A) \;=\; \frac{P(A\mid B)\,P(B)}{P(A\mid B)\,P(B) \,+\, P(A\mid B^c)\,P(B^c)}.
  \]
  $P(B)$: probabilidad \emph{a priori}. \quad $P(B \mid A)$: probabilidad \emph{a posteriori}.
\end{teorema}

\vspace{4pt}
\begin{center}
  \begin{tikzpicture}[
      every node/.style={font=\small},
      level 1/.style={sibling distance=5cm, level distance=2cm},
      level 2/.style={sibling distance=2.2cm, level distance=2cm},
      edge from parent/.style={draw, -Latex},
    ]
    \node[circle,draw,fill=sectionbg!15,minimum size=8mm] {$\Omega$}
    child { node[circle,draw,fill=defbg,draw=defborder,minimum size=8mm] {$B_1$}
        child { node[rectangle,draw,rounded corners,fill=exbg] {$A\cap B_1$}
            edge from parent node[left,font=\tiny] {$P(A|B_1)$} }
        child { node[rectangle,draw,rounded corners,fill=exbg] {$A^c\cap B_1$}
            edge from parent node[right,font=\tiny] {$1-P(A|B_1)$} }
        edge from parent node[left] {$P(B_1)$}
      }
    child { node[circle,draw,fill=notebg,draw=noteborder,minimum size=8mm] {$B_2$}
        child { node[rectangle,draw,rounded corners,fill=exbg] {$A\cap B_2$}
            edge from parent node[left,font=\tiny] {$P(A|B_2)$} }
        child { node[rectangle,draw,rounded corners,fill=exbg] {$A^c\cap B_2$}
            edge from parent node[right,font=\tiny] {$1-P(A|B_2)$} }
        edge from parent node[right] {$P(B_2)$}
      };
  \end{tikzpicture}
\end{center}
\begin{center}
  {\footnotesize\itshape \'Arbol de probabilidad condicional con partici\'on $\{B_1, B_2\}$.}
\end{center}

\begin{ejemplo}{Prueba diagn\'ostica (Bayes en acci\'on)}
  \textbf{Notaci\'on.} Sean los eventos:
  \begin{itemize}
    \item $B = $ ``la persona est\'a enferma''.
    \item $A = $ ``el test da positivo''.
  \end{itemize}
  \textbf{Datos del problema:}
  \begin{itemize}
    \item Prevalencia: $P(B) = 0.01$ \hfill (el $1\%$ de la poblaci\'on est\'a enferma).
    \item Sensibilidad: $P(A\mid B) = 0.95$ \hfill ($95\%$ de los enfermos dan positivo).
    \item Especificidad: $P(A^c\mid B^c) = 0.90$ \hfill ($90\%$ de los sanos dan negativo).
  \end{itemize}
  De la especificidad: $P(A\mid B^c) = 1 - 0.90 = 0.10$ (tasa de falsos positivos).
  \\[5pt]
  \textbf{Pregunta:} si una persona da positivo, \textquestiondown cu\'al es la probabilidad
  de que est\'e enferma? Es decir, calcular $P(B\mid A)$.
  \\[5pt]
  \textbf{Paso 1.} Calcular $P(A)$ por probabilidad total con la partici\'on $\{B, B^c\}$:
  \[
    P(A) = P(A\mid B)\,P(B) + P(A\mid B^c)\,P(B^c)
    = 0.95\cdot 0.01 + 0.10\cdot 0.99 = 0.0095 + 0.099 = 0.1085.
  \]
  \textbf{Paso 2.} Aplicar Bayes:
  \[
    P(B\mid A) = \frac{P(A\mid B)\,P(B)}{P(A)}
    = \frac{0.95\cdot 0.01}{0.1085} = \frac{0.0095}{0.1085} \approx 0.0875.
  \]
  \textbf{Interpretaci\'on.} Aunque el test tiene $95\%$ de sensibilidad, solo el $\sim 8.8\%$
  de quienes dan positivo est\'an realmente enfermos: el $\sim 91\%$ son \emph{falsos positivos}.
  Esta es la \emph{paradoja del valor predictivo positivo}: cuando una enfermedad es rara, la
  poblaci\'on sana es mucho m\'as grande y aunque su tasa de error sea baja, genera m\'as
  positivos en n\'umeros absolutos.
\end{ejemplo}

\subsecheader{2.2 \textperiodcentered\ Independencia}

\begin{definicion}{Independencia de eventos}
  $A, B \in \sigalg$ son \textbf{independientes} ($A \perp B$) si:
  \[
    P(A \cap B) = P(A)\,P(B).
  \]
  Una familia $\{A_i\}_{i\in I}$ es \textbf{mutuamente independiente} si para todo $J\subseteq I$
  finito:
  \[
    P\!\left(\bigcap_{j\in J} A_j\right) = \prod_{j \in J} P(A_j).
  \]
\end{definicion}

\begin{cuidado}
  Independencia mutua es m\'as fuerte que independencia por pares. Existen familias donde
  todos los pares son independientes pero la familia completa no lo es
  (contraejemplo: dados de Bernstein con $\Omega = \{1,2,3,4\}$, $P$ uniforme).
\end{cuidado}

% =============================================================================
\secheader{3}{Variables Aleatorias}

\subsecheader{3.1 \textperiodcentered\ Definici\'on medida-te\'orica}

\begin{definicion}{$\sigma$-\'algebra de Borel y conjuntos borelianos}
  La \textbf{$\sigma$-\'algebra de Borel} sobre $\R$, denotada $\Borel(\R)$, es la
  \emph{m\'as peque\~na} $\sigma$-\'algebra que contiene a todos los intervalos:
  \[
    \Borel(\R) \;=\; \sigma\!\left(\{(a,b] : a<b,\; a,b\in\R\}\right).
  \]
  Sus elementos se llaman \textbf{conjuntos borelianos} (o \emph{borelianos}). Por construcci\'on,
  $\Borel(\R)$ contiene a todos los:
  \begin{itemize}
    \item intervalos abiertos $(a,b)$, cerrados $[a,b]$, semi-abiertos $(a,b]$ y $[a,b)$;
    \item rayos $(-\infty, a]$, $(a, \infty)$, etc.;
    \item conjuntos abiertos y cerrados de $\R$;
    \item uniones e intersecciones \emph{numerables} de los anteriores;
    \item complementos de cualquier boreliano.
  \end{itemize}
  En la pr\'actica, todo subconjunto de $\R$ que aparece naturalmente (intervalos, puntos,
  uniones contables, etc.) es boreliano.
\end{definicion}

\begin{nota}
  \textbf{\textquestiondown Por qu\'e no usar todos los subconjuntos de $\R$?} Igual que con
  la $\sigma$-\'algebra general, en $\R$ existen conjuntos ``patol\'ogicos'' (no medibles, como
  los conjuntos de Vitali) sobre los cuales no se puede definir una medida de probabilidad
  consistente. $\Borel(\R)$ es el ``contenedor'' suficientemente rico para todos los eventos
  de inter\'es y suficientemente restrictivo para excluir los problem\'aticos.
\end{nota}

\begin{definicion}{Variable aleatoria}
  Dado $(\Omega, \sigalg, P)$, una \textbf{variable aleatoria} es una funci\'on medible
  $X : \Omega \to \R$, es decir:
  \[
    \forall\, B \in \Borel(\R): \quad X^{-1}(B) = \{\omega \in \Omega : X(\omega) \in B\} \in \sigalg.
  \]
  Equivalentemente, basta verificar que $\{\omega : X(\omega) \leq x\} \in \sigalg$ para todo $x\in\R$.
\end{definicion}

\begin{center}
  \begin{tikzpicture}[scale=0.85, >=Latex]
    \draw[thick, rounded corners=6pt, fill=thmbg!40] (-3.5,-1.8) rectangle (-0.2,1.8);
    \node[font=\small\bfseries, color=sectionbg] at (-1.85,1.5) {$\Omega$};
    \draw[fill=defbg, draw=defborder] (-1.85,0) ellipse (1.1cm and 0.9cm);
    \node[font=\tiny, color=defborder] at (-1.85,0.15) {$X^{-1}(B)$};
    \node[font=\tiny, color=defborder] at (-1.85,-0.25) {$\in\sigalg$};
    \draw[thick, -Latex] (0.5,-0.1) -- (4.5,-0.1) node[right,font=\small] {$\R$};
    \draw[very thick, color=noteborder] (1.8,-0.1) -- (3.2,-0.1);
    \node[above, font=\tiny\bfseries, color=noteborder] at (2.5,0.05) {$B \in \Borel(\R)$};
    \draw[thin, color=noteborder] (1.8,-0.05) -- (1.8,-0.18) (3.2,-0.05) -- (3.2,-0.18);
    \draw[-Latex, very thick, color=codeblue] (-0.2,0.2) to[bend left=25]
    node[above, font=\small\bfseries, color=codeblue] {$X$} (1.6,-0.1);
  \end{tikzpicture}
\end{center}
\begin{center}
  {\footnotesize\itshape $X$ lleva $\omega\in\Omega$ a $\R$; la pre-imagen de cualquier boreliano
    debe pertenecer a $\sigalg$.}
\end{center}

\subsecheader{3.2 \textperiodcentered\ Funci\'on de distribuci\'on acumulada (FDA)}

\begin{definicion}{Funci\'on de distribuci\'on acumulada (FDA)}
  La \textbf{funci\'on de distribuci\'on acumulada (FDA)} de $X$ es:
  \[
    F_X(x) = P(X \leq x), \qquad x \in \R.
  \]
\end{definicion}

\begin{proposicion}{Propiedades de la FDA}
  Toda FDA $F_X$ satisface las siguientes tres propiedades:
  \\[3pt]
  \textbf{(1) Mon\'otona no decreciente:} $\;x \leq y \;\Rightarrow\; F(x) \leq F(y)$.
  \begin{itemize}
    \item \emph{Por qu\'e:} si $x\leq y$, entonces $\{X\leq x\}\subseteq\{X\leq y\}$.
          Por monoton\'ia de $P$: $P(X\leq x)\leq P(X\leq y)$.
    \item \emph{Intuici\'on:} la FDA \emph{acumula} probabilidad de izquierda a derecha;
          al avanzar nunca puede ``perder'' masa.
  \end{itemize}
  \textbf{(2) Continua por la derecha:} $\;\lim_{t\downarrow x} F(t) = F(x)$.
  \begin{itemize}
    \item \emph{Notaci\'on.} $t\downarrow x$ significa que $t$ se acerca a $x$
          \textbf{desde la derecha} ($t>x$ y $t\to x$). Equivale a $\lim_{t\to x^+} F(t)$.
    \item \emph{Por qu\'e (demostraci\'on).} Si $t_n \downarrow x$, los eventos $\{X\leq t_n\}$
          forman una sucesi\'on decreciente cuya intersecci\'on es $\{X\leq x\}$.
          Por continuidad de $P$ para sucesiones decrecientes:
          $\lim P(X\leq t_n) = P(X\leq x) = F(x)$.
    \item \emph{Intuici\'on con saltos.} En un punto de salto $x_0$:
          \begin{itemize}
            \item Por la izquierda: $\lim_{t\uparrow x_0} F(t) = F(x_0^-)$ (valor \emph{abajo} del salto).
            \item Por la derecha: $\lim_{t\downarrow x_0} F(t) = F(x_0)$ (valor \emph{arriba} del salto).
          \end{itemize}
          La propiedad (2) dice que $F$ ``salta y se queda arriba''.
    \item \emph{Tama\~no del salto:} $F(x_0) - F(x_0^-) = P(X = x_0)$.
  \end{itemize}
  \textbf{(3) L\'imites:} $\;\lim_{x\to-\infty} F(x) = 0$ \;y\; $\lim_{x\to+\infty} F(x) = 1$.
  \begin{itemize}
    \item \emph{Por qu\'e:} cuando $x\to-\infty$, los eventos $\{X\leq x\}$ decrecen a $\emptyset$.
          Cuando $x\to+\infty$, crecen a $\Omega$. Por continuidad de $P$: $P(\emptyset)=0$, $P(\Omega)=1$.
  \end{itemize}
  \textbf{Rec\'iproco.} Toda funci\'on $F:\R\to[0,1]$ que satisfaga (1), (2) y (3) es la FDA
  de \emph{alguna} variable aleatoria. Por tanto, la FDA caracteriza completamente la
  distribuci\'on de $X$.
\end{proposicion}

\subsecheader{3.3 \textperiodcentered\ Variables discretas y continuas}

\begin{definicion}{Variable aleatoria discreta y funci\'on de masa (PMF)}
  $X$ es \textbf{discreta} si toma valores en un conjunto numerable $\{x_1, x_2, \ldots\}$.
  Su \textbf{funci\'on de masa de probabilidad (PMF)} es:
  \[
    p_X(x_k) = P(X = x_k) \geq 0, \qquad \sum_k p_X(x_k) = 1.
  \]
\end{definicion}

\begin{definicion}{Variable aleatoria continua y funci\'on de densidad (PDF)}
  $X$ es \textbf{continua} si existe una \textbf{funci\'on de densidad de probabilidad (PDF)}
  $f_X : \R \to [0,\infty)$ tal que:
  \[
    F_X(x) = \int_{-\infty}^{x} f_X(t)\,\mathrm{d}t, \qquad \int_{-\infty}^{\infty} f_X(t)\,\mathrm{d}t = 1.
  \]
  Para $X$ continua: $P(X = x) = 0$ para todo $x\in\R$.
\end{definicion}

\begin{definicion}{Soporte de una variable aleatoria}
  El \textbf{soporte} de $X$, denotado $I_X$, es el conjunto de valores que $X$ toma con
  probabilidad positiva:
  \[
    I_X \;=\;
    \begin{cases}
      \{x_k : p_X(x_k) > 0\} & \text{si } X \text{ es discreta,} \\
      \{x \in \R : f_X(x) > 0\}             & \text{si } X \text{ es continua.}
    \end{cases}
  \]
  En lo que sigue, sumas e integrales se restringen a $I_X$.
\end{definicion}

\subsecheader{3.4 \textperiodcentered\ Esperanza y varianza}

\begin{definicion}{Esperanza}
  \textit{Discreta:} $\;\E{X} = \displaystyle\sum_{x\in I_X} x\, p_X(x)$.
  \\[3pt]
  \textit{Continua:} $\;\E{X} = \displaystyle\int_{I_X} x\, f_X(x)\,\mathrm{d}x$.
  \\[3pt]
  Si $g:\R\to\R$ medible: $\;\E{g(X)} = \displaystyle\sum_{x\in I_X} g(x)\, p_X(x)$
  \;o\; $\displaystyle\int_{I_X} g(x)\, f_X(x)\,\mathrm{d}x$. \quad
  (\emph{Ley del estad\'istico inconsciente})
\end{definicion}

\begin{proposicion}{Propiedades de la esperanza}
  Sean $X, Y$ variables aleatorias con esperanza finita y $a, b, c \in \R$ constantes.
  \begin{enumerate}
    \item \textbf{Constante:} $\;\E{c} = c$.
    \item \textbf{Linealidad:} $\;\E{aX + bY} = a\,\E{X} + b\,\E{Y}$.
          \emph{Vale siempre, no requiere independencia.}
    \item \textbf{Monoton\'ia:} si $X \leq Y$ casi seguramente, entonces $\E{X} \leq \E{Y}$.
    \item \textbf{No negatividad:} si $X \geq 0$ casi seguramente, entonces $\E{X} \geq 0$.
    \item \textbf{Desigualdad triangular:} $\;|\E{X}| \leq \E{|X|}$.
    \item \textbf{Producto bajo independencia:} si $X \perp Y$, entonces $\E{XY} = \E{X}\,\E{Y}$.
  \end{enumerate}
\end{proposicion}

\begin{nota}
  \textbf{Linealidad} es la propiedad m\'as utilizada en pr\'actica: permite descomponer
  esperanzas de sumas en sumas de esperanzas \emph{aunque las variables sean dependientes}.
\end{nota}

\begin{definicion}{Varianza}
  \[
    \Var{X} = \E{(X-\E{X})^2} = \E{X^2} - (\E{X})^2.
  \]
  \textbf{Desviaci\'on est\'andar}: $\sigma_X = \sqrt{\Var{X}}$.
\end{definicion}

\begin{proposicion}{Propiedades de la varianza}
  Sean $X, Y$ variables aleatorias con varianza finita y $a, b \in \R$ constantes.
  \begin{enumerate}
    \item \textbf{No negatividad:} $\;\Var{X} \geq 0$, con $\Var{X} = 0$ si y solo si $X$ es constante c.s.
    \item \textbf{Constante:} $\;\Var{c} = 0$.
    \item \textbf{Translaci\'on (invariancia):} $\;\Var{X + b} = \Var{X}$.
    \item \textbf{Escalamiento:} $\;\Var{aX} = a^2\,\Var{X}$.
    \item \textbf{Transformaci\'on af\'in:} $\;\Var{aX + b} = a^2\,\Var{X}$.
    \item \textbf{Suma (caso general):} $\;\Var{X + Y} = \Var{X} + \Var{Y} + 2\,\Cov{X}{Y}$,
          donde $\Cov{X}{Y} = \E{XY} - \E{X}\E{Y}$.
    \item \textbf{Suma bajo independencia:} si $X \perp Y$:
          $\;\Var{X + Y} = \Var{X} + \Var{Y}$. M\'as general:
          $\Var{\sum X_i} = \sum \Var{X_i}$ si todos son independientes.
  \end{enumerate}
\end{proposicion}

\begin{nota}
  \textbf{Cuidado:} a diferencia de la esperanza, \emph{la varianza no es lineal}.
  $\Var{X+Y} \neq \Var{X} + \Var{Y}$ en general --- solo bajo independencia
  (o, m\'as d\'ebilmente, descorrelaci\'on $\Cov{X}{Y}=0$). La desviaci\'on est\'andar
  s\'i escala linealmente: $\sigma_{aX} = |a|\,\sigma_X$.
\end{nota}

% =============================================================================
\secheader{4}{Distribuciones de Probabilidad Comunes}

\subsecheader{4.1 \textperiodcentered\ Tabla resumen}

\begin{center}
  \renewcommand{\arraystretch}{1.4}
  \begin{tabular}{@{}l l l c c l@{}}
    \toprule
    \textbf{Distribuci\'on} & \textbf{Par\'ametros}    & \textbf{Soporte} & $\mathbb{E}[X]$  & $\mathrm{Var}(X)$     & \textbf{MGF}                              \\
    \midrule
    Bernoulli($p$)          & $p\in(0,1)$              & $\{0,1\}$        & $p$              & $p(1-p)$              & $1-p+pe^t$                                \\
    Binomial($n,p$)         & $n\in\Nset$, $p\in(0,1)$ & $\{0,\ldots,n\}$ & $np$             & $np(1-p)$             & $(1-p+pe^t)^n$                            \\
    Geom\'etrica($p$)       & $p\in(0,1)$              & $\Nset^+$        & $1/p$            & $(1-p)/p^2$           & $\dfrac{pe^t}{1-(1-p)e^t}$                \\[4pt]
    Poisson($\lambda$)      & $\lambda>0$              & $\Nset_0$        & $\lambda$        & $\lambda$             & $e^{\lambda(e^t-1)}$                      \\
    Uniforme($a,b$)         & $a<b$                    & $[a,b]$          & $\tfrac{a+b}{2}$ & $\tfrac{(b-a)^2}{12}$ & $\tfrac{e^{tb}-e^{ta}}{t(b-a)}$           \\
    Exponencial($\lambda$)  & $\lambda>0$              & $[0,\infty)$     & $1/\lambda$      & $1/\lambda^2$         & $\tfrac{\lambda}{\lambda-t},\; t<\lambda$ \\
    Normal($\mu,\sigma^2$)  & $\mu\in\R,\sigma>0$      & $\R$             & $\mu$            & $\sigma^2$            & $e^{\mu t+\sigma^2 t^2/2}$                \\
    \bottomrule
  \end{tabular}
\end{center}

\vspace{6pt}
\begin{nota}
  \textbf{Mapa de decisi\'on r\'apida.}
  \begin{itemize}\setlength{\itemsep}{1pt}
    \item \textquestiondown Cuento o mido? \;\emph{Cuento} $\to$ discreta. \;\emph{Mido tiempo/distancia/cantidad continua} $\to$ continua.
    \item \textquestiondown N\'umero \emph{fijo} de intentos $n$? \;S\'i $\to$ Binomial. \;No, hasta el primer \'exito $\to$ Geom\'etrica.
          \;No, eventos a tasa por unidad $\to$ Poisson.
    \item \textquestiondown Tiempo entre eventos? \;$\to$ Exponencial.
    \item \textquestiondown Sin preferencia en un rango? \;$\to$ Uniforme.
    \item \textquestiondown Suma de muchos efectos? \;$\to$ Normal (por TCL).
  \end{itemize}
\end{nota}

\subsecheader{4.2 \textperiodcentered\ Distribuci\'on Bernoulli}

\begin{definicion}{Distribuci\'on Bernoulli}
  $X \sim \mathrm{Bernoulli}(p)$ con $p\in(0,1)$ modela un \'unico ensayo con dos resultados.
  Soporte $I_X = \{0,1\}$, PMF:
  \[
    p_X(k) = p^k (1-p)^{1-k}, \quad k\in I_X.
  \]
  Equivalentemente: $P(X=1)=p$, $P(X=0)=1-p$.
  \textbf{Momentos:} $\E{X} = p$, \quad $\Var{X} = p(1-p)$.
\end{definicion}

\begin{demostracion}
  Verificamos que la PMF suma $1$:
  \[
    \sum_{k=0}^{1} p^k(1-p)^{1-k} = (1-p) + p = 1.
  \]
\end{demostracion}

\begin{nota}
  \textbf{\textquestiondown Cu\'ando usar?} Un solo experimento con dos resultados mutuamente
  excluyentes. Ladrillo elemental: la Binomial y la Geom\'etrica se construyen como sucesiones
  de Bernoullis i.i.d.\\
  \textbf{Ejemplo:} lanzar una moneda; pregunta s\'i/no en encuesta.
\end{nota}

\begin{center}
  \begin{tikzpicture}
    \begin{axis}[
        width=11cm, height=4.5cm,
        xlabel={$k$}, ylabel={$P(X=k)$},
        xtick={0,1}, ytick={0,0.2,0.4,0.6,0.8},
        ymin=0, ymax=0.85, xmin=-0.5, xmax=1.5,
        ybar, bar width=20pt,
        grid=major, grid style={dotted, gray!40},
        tick label style={font=\small},
        label style={font=\small},
        legend style={at={(0.98,0.95)},anchor=north east,font=\small},
      ]
      \addplot[fill=defborder!70, draw=defborder, bar shift=-11pt] coordinates {(0,0.7) (1,0.3)};
      \addlegendentry{$p=0.3$}
      \addplot[fill=noteborder!70, draw=noteborder, bar shift=11pt] coordinates {(0,0.5) (1,0.5)};
      \addlegendentry{$p=0.5$}
    \end{axis}
  \end{tikzpicture}
\end{center}
\begin{center}
  {\footnotesize\itshape PMF de Bernoulli para $p=0.3$ (sesgada) y $p=0.5$ (justa).}
\end{center}

\subsecheader{4.3 \textperiodcentered\ Distribuci\'on Binomial}

\begin{definicion}{Distribuci\'on Binomial}
  $X \sim \mathrm{Bin}(n, p)$ con $n\in\Nset$, $p\in(0,1)$ cuenta el n\'umero de \'exitos en $n$
  ensayos Bernoulli$(p)$ \emph{independientes}. Soporte $I_X = \{0,1,\ldots,n\}$, PMF:
  \[
    p_X(k) = \binom{n}{k} p^k (1-p)^{n-k}, \quad k\in I_X.
  \]
  Equivalentemente, $X = \sum_{i=1}^{n} Y_i$ con $Y_i \sim \mathrm{Bernoulli}(p)$ i.i.d.
  \textbf{Momentos:} $\E{X} = np$, \quad $\Var{X} = np(1-p)$.
\end{definicion}

\begin{demostracion}
  Por el \emph{teorema del binomio} aplicado a $(p + (1-p))^n$:
  \[
    \sum_{k=0}^{n} \binom{n}{k} p^k (1-p)^{n-k}
    = \big(p + (1-p)\big)^n = 1^n = 1.
  \]
\end{demostracion}

\begin{nota}
  \textbf{\textquestiondown Cu\'ando usar?} Cuando se realizan $n$ ensayos Bernoulli
  independientes con la \emph{misma} $p$ y se cuenta el total de \'exitos.\\
  \textbf{Ejemplo:} \# caras en $10$ lanzamientos; \# defectuosas en una muestra de $100$ piezas.
\end{nota}

\begin{center}
  \begin{tikzpicture}
    \begin{axis}[
        width=13cm, height=5.5cm,
        xlabel={$k$}, ylabel={$P(X=k)$},
        xtick={0,1,2,3,4,5,6,7,8,9,10},
        ymin=0, ymax=0.32,
        legend style={at={(0.98,0.95)},anchor=north east,font=\small},
        tick label style={font=\small},
        label style={font=\small},
        grid=major, grid style={dotted, gray!40},
        ybar, bar width=5pt,
      ]
      \addplot[fill=defborder!70, draw=defborder, bar shift=-3pt] coordinates {
          (0,0.0282)(1,0.1211)(2,0.2335)(3,0.2668)(4,0.2001)
          (5,0.1029)(6,0.0368)(7,0.0090)(8,0.0014)(9,0.0001)(10,0)
        };
      \addlegendentry{$\mathrm{Bin}(10,\,0.3)$}
      \addplot[fill=noteborder!70, draw=noteborder, bar shift=3pt] coordinates {
          (0,0.0010)(1,0.0098)(2,0.0439)(3,0.1172)(4,0.2051)
          (5,0.2461)(6,0.2051)(7,0.1172)(8,0.0439)(9,0.0098)(10,0.0010)
        };
      \addlegendentry{$\mathrm{Bin}(10,\,0.5)$}
    \end{axis}
  \end{tikzpicture}
\end{center}
\begin{center}
  {\footnotesize\itshape PMF de Binomial$(n=10)$. Con $p=0.5$ es sim\'etrica;
    con $p\neq 0.5$ se sesga hacia $np$.}
\end{center}

\subsecheader{4.4 \textperiodcentered\ Distribuci\'on Geom\'etrica}

\begin{definicion}{Distribuci\'on Geom\'etrica}
  $X \sim \mathrm{Geom}(p)$ con $p\in(0,1)$ cuenta el n\'umero de ensayos Bernoulli$(p)$ i.i.d.\
  \emph{hasta el primer \'exito} (incluido). Soporte $I_X = \{1,2,3,\ldots\}$, PMF:
  \[
    p_X(k) = p\,(1-p)^{k-1}, \quad k\in I_X.
  \]
  \textbf{Momentos:} $\E{X} = 1/p$, \quad $\Var{X} = (1-p)/p^2$.
\end{definicion}

\begin{demostracion}
  Reindexando $j = k-1$ y usando la \emph{serie geom\'etrica}:
  \[
    \sum_{k=1}^{\infty} p\,(1-p)^{k-1}
    = p\sum_{j=0}^{\infty} (1-p)^j
    = p \cdot \frac{1}{1-(1-p)} = p \cdot \frac{1}{p} = 1.
  \]
\end{demostracion}

\begin{nota}
  \textbf{\textquestiondown Cu\'ando usar?} Cuando se repiten ensayos Bernoulli y se cuenta
  cu\'antos ensayos hacen falta hasta el primer \'exito (no hay tope $n$).\\
  \textbf{Ejemplo:} \# lanzamientos hasta la primera cara; \# productos hasta el primero defectuoso.
\end{nota}

\begin{proposicion}{P\'erdida de memoria (caso discreto)}
  $X \sim \mathrm{Geom}(p)$ si y s\'olo si:
  \[
    P(X > m+n \mid X > m) = P(X > n) \quad \forall\, m, n \in \Nset_0.
  \]
  Es la \emph{\'unica} distribuci\'on discreta con esta propiedad (an\'aloga discreta de la Exponencial).
\end{proposicion}

\begin{demostracion}
  \textbf{($\Rightarrow$)} La \emph{funci\'on de supervivencia} es $P(X > k) = (1-p)^k$. Como
  $m+n \geq m$, $\{X>m+n\} \subseteq \{X>m\}$, luego:
  \[
    P(X > m+n \mid X > m) = \frac{(1-p)^{m+n}}{(1-p)^m} = (1-p)^n = P(X > n). \;\checkmark
  \]
  \textbf{($\Leftarrow$)} Sea $G(k) = P(X > k)$ con $G(0) = 1$. La propiedad equivale a
  $G(m+n) = G(m)\,G(n)$. Por inducci\'on, $G(k) = G(1)^k = q^k$. La PMF resulta
  $p_X(k) = q^{k-1}(1-q) = (1-p)^{k-1}\, p$ con $p = 1-q$.
\end{demostracion}

\begin{ejemplo}{P\'erdida de memoria en acci\'on (Geom\'etrica)}
  Cada pieza es defectuosa con probabilidad $p = 0.1$. Sea $X$ = \# de piezas inspeccionadas
  hasta la primera defectuosa. Entonces $X \sim \mathrm{Geom}(0.1)$ con $\E{X} = 10$.
  \\[5pt]
  \textbf{Pregunta:} ya inspeccionaste $5$ piezas y todas resultaron buenas.
  \textquestiondown Cu\'al es la probabilidad de inspeccionar al menos $3$ piezas m\'as
  sin encontrar defectuosa?
  \\[5pt]
  \textbf{Soluci\'on:} $P(X > 8 \mid X > 5) = P(X > 3) = (0.9)^3 = 0.729$.
  \\[5pt]
  \textbf{Interpretaci\'on.} Las $5$ inspecciones previas no aportan informaci\'on:
  la probabilidad de inspeccionar otras $3$ sin defectuosa es la misma que empezar de cero.
\end{ejemplo}

\begin{center}
  \begin{tikzpicture}
    \begin{axis}[
        width=13cm, height=5.5cm,
        xlabel={$k$}, ylabel={$P(X=k)$},
        xtick={1,2,3,4,5,6,7,8,9,10,11,12},
        ymin=0, ymax=0.55,
        legend style={at={(0.98,0.95)},anchor=north east,font=\small},
        tick label style={font=\small},
        label style={font=\small},
        grid=major, grid style={dotted, gray!40},
        ybar, bar width=5pt,
      ]
      \addplot[fill=defborder!70, draw=defborder, bar shift=-3pt] coordinates {
          (1,0.3)(2,0.21)(3,0.147)(4,0.1029)(5,0.0720)
          (6,0.0504)(7,0.0353)(8,0.0247)(9,0.0173)(10,0.0121)(11,0.0085)(12,0.0059)
        };
      \addlegendentry{$p=0.3$}
      \addplot[fill=noteborder!70, draw=noteborder, bar shift=3pt] coordinates {
          (1,0.5)(2,0.25)(3,0.125)(4,0.0625)(5,0.0313)
          (6,0.0156)(7,0.0078)(8,0.0039)(9,0.0020)(10,0.0010)(11,0.0005)(12,0.0002)
        };
      \addlegendentry{$p=0.5$}
    \end{axis}
  \end{tikzpicture}
\end{center}
\begin{center}
  {\footnotesize\itshape PMF de Geom\'etrica: decrece geom\'etricamente.}
\end{center}

\subsecheader{4.5 \textperiodcentered\ Distribuci\'on Poisson}

\begin{definicion}{Distribuci\'on Poisson}
  $X \sim \mathrm{Poisson}(\lambda)$ con $\lambda > 0$ cuenta eventos raros que ocurren a
  tasa constante en un intervalo. Soporte $I_X = \Nset_0$, PMF:
  \[
    p_X(k) = \frac{\lambda^k\, e^{-\lambda}}{k!}, \quad k \in I_X.
  \]
  \textbf{Momentos:} $\E{X} = \lambda$, \quad $\Var{X} = \lambda$ (\emph{equidispersi\'on}).
\end{definicion}

\begin{demostracion}
  Usando la \emph{serie de Taylor} de la exponencial $e^\lambda = \sum_{k=0}^{\infty}\lambda^k/k!$:
  \[
    \sum_{k=0}^{\infty} \frac{\lambda^k e^{-\lambda}}{k!}
    = e^{-\lambda} \sum_{k=0}^{\infty} \frac{\lambda^k}{k!}
    = e^{-\lambda} \cdot e^{\lambda} = 1.
  \]
\end{demostracion}

\begin{nota}
  \textbf{\textquestiondown Cu\'ando usar?} Cuando se cuentan eventos que ocurren
  \emph{independientemente} a una tasa constante en tiempo o espacio. L\'imite de
  $\mathrm{Bin}(n,\lambda/n)$ cuando $n\to\infty$.\\
  \textbf{Ejemplo:} \# llamadas/hora a un call center; \# decaimientos radiactivos/seg.
\end{nota}

\begin{center}
  \begin{tikzpicture}
    \begin{axis}[
        width=13cm, height=6cm,
        xlabel={$k$}, ylabel={$P(X=k)$},
        xtick={0,1,2,3,4,5,6,7,8,9,10,11,12,13,14,15,16},
        ymin=0, ymax=0.42,
        legend style={at={(0.98,0.95)},anchor=north east,font=\small},
        tick label style={font=\small},
        label style={font=\small},
        grid=major, grid style={dotted, gray!40},
        ybar, bar width=4pt,
      ]
      \addplot[fill=sectionbg!70, draw=sectionbg, bar shift=-4pt] coordinates {
          (0,0.3679)(1,0.3679)(2,0.1839)(3,0.0613)(4,0.0153)(5,0.0031)(6,0.0005)
        };
      \addlegendentry{$\lambda=1$}
      \addplot[fill=defborder!70, draw=defborder, bar shift=0pt] coordinates {
          (0,0.0498)(1,0.1494)(2,0.2240)(3,0.2240)(4,0.1680)(5,0.1008)(6,0.0504)
          (7,0.0216)(8,0.0081)(9,0.0027)(10,0.0008)
        };
      \addlegendentry{$\lambda=3$}
      \addplot[fill=noteborder!70, draw=noteborder, bar shift=4pt] coordinates {
          (0,0.0009)(1,0.0064)(2,0.0223)(3,0.0521)(4,0.0912)(5,0.1277)(6,0.1490)
          (7,0.1490)(8,0.1304)(9,0.1014)(10,0.0710)(11,0.0452)(12,0.0264)(13,0.0142)
          (14,0.0071)(15,0.0033)
        };
      \addlegendentry{$\lambda=7$}
    \end{axis}
  \end{tikzpicture}
\end{center}
\begin{center}
  {\footnotesize\itshape PMF de Poisson para $\lambda\in\{1,3,7\}$.}
\end{center}

\subsecheader{4.6 \textperiodcentered\ Distribuci\'on Uniforme}

\begin{definicion}{Distribuci\'on Uniforme}
  $X \sim \mathrm{Unif}(a,b)$ con $a<b$ tiene soporte $I_X = [a,b]$ y densidad constante:
  \[
    f_X(x) = \frac{1}{b-a}, \quad x\in I_X, \qquad
    F_X(x) = \frac{x-a}{b-a}.
  \]
  \textbf{Momentos:} $\E{X} = (a+b)/2$, \quad $\Var{X} = (b-a)^2/12$.
\end{definicion}

\begin{demostracion}
  $\int_a^b \frac{1}{b-a}\,\mathrm{d}x = \frac{b-a}{b-a} = 1.$
\end{demostracion}

\begin{nota}
  \textbf{\textquestiondown Cu\'ando usar?} Cuando todo punto del intervalo $[a,b]$ es igualmente
  probable.\\
  \textbf{Aplicaci\'on cr\'itica:} la $\mathrm{Unif}(0,1)$ es la base computacional para generar
  cualquier distribuci\'on por \emph{transformada inversa}: $F^{-1}(U) \sim F$.
\end{nota}

\begin{center}
  \begin{tikzpicture}
    \begin{axis}[
        width=12cm, height=4.5cm,
        xlabel={$x$}, ylabel={$f(x)$},
        xtick={0,1,2,3,4,5,6},
        ymin=0, ymax=1.2, xmin=-0.5, xmax=6.5,
        grid=major, grid style={dotted, gray!40},
        tick label style={font=\small},
        label style={font=\small},
        legend style={at={(0.98,0.95)},anchor=north east,font=\small},
        axis lines=left,
      ]
      \addplot[thick, color=defborder, fill=defbg, opacity=0.7]
      coordinates {(0,0)(0,1)(1,1)(1,0)} \closedcycle;
      \addlegendentry{$\mathrm{Unif}(0,1)$}
      \addplot[thick, color=noteborder, fill=notebg, opacity=0.7]
      coordinates {(2,0)(2,0.25)(6,0.25)(6,0)} \closedcycle;
      \addlegendentry{$\mathrm{Unif}(2,6)$}
    \end{axis}
  \end{tikzpicture}
\end{center}

\subsecheader{4.7 \textperiodcentered\ Distribuci\'on Exponencial}

\begin{definicion}{Distribuci\'on Exponencial}
  $X\sim\mathrm{Exp}(\lambda)$ con $\lambda>0$ tiene soporte $I_X = [0,\infty)$ y:
  \[
    f_X(x) = \lambda e^{-\lambda x},\quad x \in I_X, \qquad F_X(x) = 1 - e^{-\lambda x}.
  \]
  \textbf{Momentos:} $\E{X} = 1/\lambda$, \quad $\Var{X} = 1/\lambda^2$.
\end{definicion}

\begin{demostracion}
  $\int_0^\infty \lambda e^{-\lambda x}\,\mathrm{d}x = \left[-e^{-\lambda x}\right]_0^\infty = 1.$
\end{demostracion}

\begin{nota}
  \textbf{\textquestiondown Cu\'ando usar?} Tiempo \emph{continuo} de espera entre eventos
  sucesivos en un proceso de Poisson de tasa $\lambda$. Tambi\'en para vidas \'utiles
  \emph{sin envejecimiento}.\\
  \textbf{Ejemplo:} tiempo entre llegadas de clientes; tiempo hasta la falla de un componente.
\end{nota}

\begin{proposicion}{Propiedad de p\'erdida de memoria}
  $X \sim \mathrm{Exp}(\lambda)$ si y s\'olo si:
  \[
    P(X > s+t \mid X > s) = P(X > t) \quad \forall\, s,t \geq 0.
  \]
  Es la \emph{\'unica} distribuci\'on continua con esta propiedad.
\end{proposicion}

\begin{demostracion}
  \textbf{($\Rightarrow$)} La supervivencia es $\overline{F}_X(x) = e^{-\lambda x}$. Como
  $\{X>s+t\} \subseteq \{X>s\}$:
  \[
    P(X > s+t \mid X > s) = \frac{e^{-\lambda(s+t)}}{e^{-\lambda s}} = e^{-\lambda t} = P(X > t). \;\checkmark
  \]
  \textbf{($\Leftarrow$)} Sea $G(x) = P(X > x)$. La propiedad da $G(s+t) = G(s)\,G(t)$ (ecuaci\'on
  de Cauchy multiplicativa). Bajo continuidad de $G$, las \'unicas soluciones son $G(x) = e^{-\lambda x}$.
\end{demostracion}

\begin{ejemplo}{P\'erdida de memoria en acci\'on}
  Call center con $\lambda = 4$ llamadas/hora ($1/\lambda = 15$ min). Llevas $10$ min sin
  llamadas. \textquestiondown Probabilidad de esperar al menos $5$ min m\'as?
  \[
    P(X > s+t \mid X > s) = P(X > t) = e^{-4/12} = e^{-1/3} \approx 0.7165.
  \]
  El tiempo ya esperado \emph{no aporta informaci\'on}.
\end{ejemplo}

\begin{center}
  \begin{tikzpicture}
    \begin{axis}[
        width=13cm, height=5.5cm,
        xlabel={$x$}, ylabel={$f(x)$},
        xtick={0,1,2,3,4,5},
        ymin=0, ymax=2.1, xmin=0, xmax=5,
        domain=0:5, samples=100,
        legend style={at={(0.98,0.95)},anchor=north east,font=\small},
        tick label style={font=\small},
        label style={font=\small},
        grid=major, grid style={dotted, gray!40},
        axis lines=left,
      ]
      \addplot[thick, color=sectionbg] {2*exp(-2*x)};
      \addlegendentry{$\lambda=2$}
      \addplot[thick, color=defborder, dashed] {1*exp(-1*x)};
      \addlegendentry{$\lambda=1$}
      \addplot[thick, color=noteborder, dash dot] {0.5*exp(-0.5*x)};
      \addlegendentry{$\lambda=0.5$}
    \end{axis}
  \end{tikzpicture}
\end{center}

\subsecheader{4.8 \textperiodcentered\ Distribuci\'on Normal}

\begin{definicion}{Distribuci\'on Normal}
  $X \sim \mathcal{N}(\mu, \sigma^2)$ tiene soporte $I_X = \R$ y densidad:
  \[
    f_X(x) = \frac{1}{\sigma\sqrt{2\pi}}\exp\!\left(-\frac{(x-\mu)^2}{2\sigma^2}\right), \quad x\in\R.
  \]
  \textbf{Normal est\'andar} $Z \sim \mathcal{N}(0,1)$: $\phi(z) = \frac{1}{\sqrt{2\pi}}e^{-z^2/2}$,
  FDA: $\Phi(z)$.\\
  \textbf{Estandarizaci\'on:} $Z = (X-\mu)/\sigma \sim \mathcal{N}(0,1)$.
  \textbf{Momentos:} $\E{X} = \mu$, \quad $\Var{X} = \sigma^2$.
\end{definicion}

\begin{demostracion}
  Cambio de variable $u = (x-\mu)/\sigma$, $\mathrm{d}u = \mathrm{d}x/\sigma$:
  \[
    \int_{-\infty}^{\infty} \frac{1}{\sigma\sqrt{2\pi}} e^{-(x-\mu)^2/(2\sigma^2)}\,\mathrm{d}x
    = \frac{1}{\sqrt{2\pi}} \int_{-\infty}^{\infty} e^{-u^2/2}\,\mathrm{d}u = 1.
  \]
  Donde se us\'o la \emph{integral de Gauss} $\int_{-\infty}^{\infty} e^{-u^2/2}\,\mathrm{d}u = \sqrt{2\pi}$
  (cuya prueba pasa a coordenadas polares).
\end{demostracion}

\begin{nota}
  \textbf{\textquestiondown Cu\'ando usar?} La variable resulta de la \emph{suma de muchos
    efectos peque\~nos e independientes} (Teorema Central del L\'imite).\\
  \textbf{Propiedad clave:} cerrada bajo combinaciones lineales: si
  $X_i \sim \mathcal{N}(\mu_i,\sigma_i^2)$ independientes,
  $\sum a_i X_i \sim \mathcal{N}(\sum a_i\mu_i, \sum a_i^2\sigma_i^2)$.
\end{nota}

\begin{center}
  \begin{tikzpicture}
    \begin{axis}[
        width=13cm, height=5.5cm,
        domain=-4:4, samples=200,
        axis lines=left,
        xlabel={$x$}, ylabel={$f(x)$},
        xtick={-3,-2,-1,0,1,2,3},
        ytick={0,0.1,0.2,0.3,0.4},
        ymin=0, ymax=0.85,
        legend style={at={(0.98,0.95)},anchor=north east,font=\small},
        tick label style={font=\small},
        label style={font=\small},
        grid=major, grid style={dotted, gray!40},
      ]
      \addplot[thick, color=sectionbg, domain=-4:4] {exp(-x^2/2)/sqrt(2*pi)};
      \addlegendentry{$\mathcal{N}(0,1)$}
      \addplot[thick, color=defborder, dashed, domain=-4:4]
      {exp(-(x-1.5)^2/(2*0.5^2))/(0.5*sqrt(2*pi))};
      \addlegendentry{$\mathcal{N}(1.5,\,0.25)$}
      \addplot[thick, color=noteborder, dash dot, domain=-4:4]
      {exp(-x^2/(2*4))/(2*sqrt(2*pi))};
      \addlegendentry{$\mathcal{N}(0,\,4)$}
      \addplot[fill=sectionbg!15, draw=none, domain=-1:1] {exp(-x^2/2)/sqrt(2*pi)} \closedcycle;
      \node[font=\tiny, color=sectionbg] at (axis cs:0,0.18) {$68\%$};
    \end{axis}
  \end{tikzpicture}
\end{center}

\subsecheader{4.9 \textperiodcentered\ Relaciones entre distribuciones}

\begin{tcolorbox}[
    enhanced,
    colback=sectionbg!8, colframe=sectionbg,
    boxrule=0.8pt, arc=4pt,
    left=8pt, right=8pt, top=6pt, bottom=6pt,
    coltitle=white, colbacktitle=sectionbg,
    titlerule=0pt,
    fonttitle=\small\bfseries\sffamily,
    title={Conexiones entre las distribuciones de la Secci\'on 4},
  ]
  \small
  \begin{itemize}
    \item \textbf{Bernoulli $\to$ Binomial:} $\mathrm{Bin}(n,p)$ es la suma de $n$ Bernoulli$(p)$
          i.i.d. Caso particular: $\mathrm{Bin}(1,p) = \mathrm{Bernoulli}(p)$.
    \item \textbf{Bernoulli $\to$ Geom\'etrica:} $\mathrm{Geom}(p)$ cuenta cu\'antos ensayos
          Bernoulli$(p)$ hacen falta para el primer \'exito.
    \item \textbf{Binomial $\to$ Poisson:} $\mathrm{Bin}(n,\lambda/n) \xrightarrow{n\to\infty} \mathrm{Poisson}(\lambda)$.
    \item \textbf{Poisson $\leftrightarrow$ Exponencial:} en un proceso de Poisson, los tiempos
          entre llegadas son Exp$(\lambda)$ i.i.d.
    \item \textbf{Geom\'etrica $\leftrightarrow$ Exponencial:} an\'aloga discreta y continua.
    \item \textbf{Suma de Poissons:} Poisson$(\lambda_1)$ + Poisson$(\lambda_2)$ = Poisson$(\lambda_1+\lambda_2)$.
    \item \textbf{Suma de Normales:} cerrada bajo combinaciones lineales independientes.
    \item \textbf{TCL $\to$ Normal:} $\mathrm{Bin}(n,p) \approx \mathcal{N}(np, np(1-p))$ y
          $\mathrm{Poisson}(\lambda) \approx \mathcal{N}(\lambda, \lambda)$ para par\'ametros grandes.
    \item \textbf{Uniforme $\to$ cualquier distribuci\'on:} si $U\sim\mathrm{Unif}(0,1)$ y $F$
          es una FDA, entonces $F^{-1}(U) \sim F$.
  \end{itemize}
\end{tcolorbox}

% =============================================================================
\secheader{5}{Vectores Aleatorios y Distribuciones Conjuntas}

\subsecheader{5.1 \textperiodcentered\ Vectores aleatorios}

\begin{definicion}{Vector aleatorio}
  Un \textbf{vector aleatorio} de dimensi\'on $n$ es una tupla
  $\mathbf{X} = (X_1, X_2, \ldots, X_n)$ donde cada $X_i$ es una variable aleatoria definida
  sobre el mismo espacio de probabilidad $(\Omega, \sigalg, P)$. Equivalentemente,
  $\mathbf{X}: \Omega \to \R^n$ medible.
\end{definicion}

\begin{nota}
  Un vector aleatorio modela varias mediciones \emph{simult\'aneas} sobre el mismo experimento.\\
  \textbf{Ejemplos:} (altura, peso) de una persona; resultados de tres lanzamientos de un dado;
  (tiempo de espera, n\'umero de clientes) en una cola.\\
  Las componentes pueden \emph{depender} entre s\'i: estudiar cada $X_i$ por separado no captura
  toda la informaci\'on del vector.
\end{nota}

\vspace{4pt}
\begin{center}
  \begin{tikzpicture}[scale=0.85, >=Latex]
    \draw[thick, rounded corners=6pt, fill=thmbg!40] (-3.5,-1.8) rectangle (-0.2,1.8);
    \node[font=\small\bfseries, color=sectionbg] at (-1.85,1.5) {$\Omega$};
    \fill[color=sectionbg] (-2.2,0.4) circle (1.5pt);
    \node[font=\tiny, right, color=sectionbg] at (-2.1,0.4) {$\omega$};
    \draw[-Latex, thick] (1,0) -- (5,0) node[right,font=\small] {$X$};
    \draw[-Latex, thick] (3,-1.5) -- (3,2) node[above,font=\small] {$Y$};
    \fill[color=defborder] (3.8,1.0) circle (2.5pt);
    \node[above right, font=\tiny\bfseries, color=defborder] at (3.85,1.05) {$(X(\omega), Y(\omega))$};
    \draw[-Latex, very thick, color=codeblue] (-0.2,0.4) to[bend left=20]
    node[above, font=\small\bfseries, color=codeblue] {$\mathbf{X} = (X, Y)$} (3.8,1.0);
    \node[font=\small, color=sectionbg] at (4,-1.3) {$\R^2$};
  \end{tikzpicture}
\end{center}
\begin{center}
  {\footnotesize\itshape Un vector aleatorio $\mathbf{X} = (X, Y)$ asigna a cada $\omega \in \Omega$
    un par $(X(\omega), Y(\omega)) \in \R^2$.}
\end{center}

\subsecheader{5.2 \textperiodcentered\ Caso general $n$-variado}

\begin{definicion}{FDA conjunta (caso general)}
  Para un vector aleatorio $(X_1, \ldots, X_n)$, la \textbf{FDA conjunta} es:
  \[
    F_{X_1,\ldots,X_n}(x_1, \ldots, x_n) = P(X_1 \leq x_1,\, X_2 \leq x_2,\, \ldots,\, X_n \leq x_n),
    \quad (x_1,\ldots,x_n) \in \R^n.
  \]
\end{definicion}

\begin{definicion}{FDA marginal (caso general)}
  Sea $J = \{i_1, \ldots, i_k\} \subseteq \{1, \ldots, n\}$ un subconjunto de \'indices.
  La \textbf{FDA marginal} del subvector $(X_{i_1}, \ldots, X_{i_k})$ se obtiene tomando
  l\'imite $\to +\infty$ en \emph{todas las variables que no est\'an en} $J$:
  \[
    F_{X_{i_1},\ldots,X_{i_k}}(x_{i_1}, \ldots, x_{i_k})
    = \lim_{x_j \to \infty,\; j\notin J} F_{X_1,\ldots,X_n}(x_1, \ldots, x_n).
  \]
\end{definicion}

\begin{nota}
  \textbf{Regla general para marginales:} para obtener la marginal del subvector $\mathbf{X}_J$,
  \emph{elimina} las variables fuera de $J$:
  \begin{itemize}
    \item En la \textbf{FDA}: hacerlas $\to +\infty$.
    \item En la \textbf{PMF}: \emph{sumar} sobre todo su soporte (caso discreto).
    \item En la \textbf{PDF}: \emph{integrar} sobre todo su soporte (caso continuo).
  \end{itemize}
  La conjunta determina un\'ivocamente \emph{todas} las marginales (univariadas, bivariadas, etc.),
  \emph{pero NO al rev\'es}.
\end{nota}

\begin{ejemplo}{Marginales de un vector trivariado continuo}
  \textbf{Planteamiento.} Sea $(X, Y, Z)$ un vector aleatorio continuo con PDF conjunta:
  \[
    f_{X,Y,Z}(x, y, z) =
    \begin{cases}
      8xyz & \text{si } (x,y,z) \in [0,1]^3, \\
      0    & \text{en otro caso}.
    \end{cases}
  \]
  Vamos a calcular las marginales en distintas dimensiones y verificar consistencia.
  \\[5pt]
  \textbf{(1) Verificaci\'on de normalizaci\'on.} Como las variables se separan en producto,
  podemos integrar cada una de forma independiente:
  \[
    \iiint_{[0,1]^3} 8xyz\,\mathrm{d}x\,\mathrm{d}y\,\mathrm{d}z
    = 8 \cdot \int_0^1 x\,\mathrm{d}x \cdot \int_0^1 y\,\mathrm{d}y \cdot \int_0^1 z\,\mathrm{d}z
    = 8 \cdot \frac{1}{2} \cdot \frac{1}{2} \cdot \frac{1}{2} = 1. \;\checkmark
  \]
  Por tanto $f_{X,Y,Z}$ es una PDF v\'alida.
  \\[5pt]
  \textbf{(2) Marginal bivariada de $(X, Y)$:} ``eliminamos'' la variable $z$ integrando
  sobre todo su soporte $[0,1]$:
  \[
    f_{X,Y}(x, y) = \int_0^1 8xyz\,\mathrm{d}z
    = 8xy \cdot \left[\frac{z^2}{2}\right]_0^1
    = 8xy \cdot \frac{1}{2} = 4xy, \quad (x,y) \in [0,1]^2.
  \]
  \emph{Justificaci\'on:} la marginal $f_{X,Y}$ describe la densidad del par $(X,Y)$ sin
  importar qu\'e valor tom\'o $Z$; por eso ``acumulamos'' (integramos) sobre todos los $z$
  posibles.
  \\[5pt]
  \textbf{(3) Marginal univariada de $X$:} podemos obtenerla integrando $f_{X,Y}$ sobre $y$:
  \[
    f_X(x) = \int_0^1 4xy\,\mathrm{d}y
    = 4x \cdot \left[\frac{y^2}{2}\right]_0^1 = 4x \cdot \frac{1}{2} = 2x,
    \quad x \in [0,1].
  \]
  Por simetr\'ia, $f_Y(y) = 2y$ y $f_Z(z) = 2z$.
  \\[5pt]
  \textbf{(4) Consistencia con la integral directa.} La marginal de $X$ tambi\'en se obtiene
  integrando la conjunta sobre $y$ y $z$ a la vez (sin pasar por $f_{X,Y}$):
  \[
    f_X(x) = \int_0^1\!\!\int_0^1 8xyz\,\mathrm{d}y\,\mathrm{d}z
    = 8x \cdot \int_0^1 y\,\mathrm{d}y \cdot \int_0^1 z\,\mathrm{d}z
    = 8x \cdot \frac{1}{2} \cdot \frac{1}{2} = 2x. \;\checkmark
  \]
  \emph{Justificaci\'on:} obtener $f_X$ de la conjunta $f_{X,Y,Z}$ o pasar antes por la
  bivariada $f_{X,Y}$ \emph{da el mismo resultado}, porque la integraci\'on de Fubini-Tonelli
  permite hacer las integrales en cualquier orden.
  \\[5pt]
  \textbf{(5) Verificaci\'on de que las marginales son v\'alidas:}
  $\int_0^1 2x\,\mathrm{d}x = [x^2]_0^1 = 1$. \;\checkmark
  \\[5pt]
  \textbf{(6) \textquestiondown Son $X, Y, Z$ independientes?} Verificamos si la conjunta
  factoriza como producto de marginales:
  \[
    f_X(x)\,f_Y(y)\,f_Z(z) = 2x \cdot 2y \cdot 2z = 8xyz = f_{X,Y,Z}(x,y,z). \;\checkmark
  \]
  S\'i factoriza, luego $X, Y, Z$ son \emph{mutuamente independientes}, cada una con densidad
  $2u$ sobre $[0,1]$ (que favorece valores cercanos a $1$).
  \\[5pt]
  \textbf{Lecci\'on clave.} Para extraer cualquier marginal:
  \begin{itemize}
    \item identifica las variables \emph{a conservar} (subconjunto $J$);
    \item integra la conjunta sobre el resto, usando los soportes correspondientes.
  \end{itemize}
  El resultado no depende del orden en que integres las variables eliminadas.
\end{ejemplo}

\subsecheader{5.3 \textperiodcentered\ Caso bivariado ($n=2$)}

Particularizamos a $n=2$ por su importancia y por permitir visualizaci\'on directa.

\begin{definicion}{FDA conjunta y marginales (caso bivariado)}
  Para un vector $(X, Y)$:
  \[
    F_{X,Y}(x, y) = P(X \leq x,\, Y \leq y), \qquad (x,y) \in \R^2.
  \]
  Las \textbf{FDA marginales} se obtienen tomando l\'imite en la otra variable:
  \[
    F_X(x) = \lim_{y \to \infty} F_{X,Y}(x, y), \qquad
    F_Y(y) = \lim_{x \to \infty} F_{X,Y}(x, y).
  \]
\end{definicion}

\begin{proposicion}{Propiedades de la FDA bivariada}
  $F_{X,Y}: \R^2 \to [0,1]$ satisface:
  \begin{enumerate}
    \item \textbf{Mon\'otona no decreciente} en cada variable.
    \item \textbf{Continua por la derecha} en cada variable.
    \item \textbf{L\'imites:} $F_{X,Y}(x,y) \to 0$ si $x\to-\infty$ o $y\to-\infty$;
          $\to 1$ si $x,y\to+\infty$ simult\'aneamente.
    \item \textbf{Probabilidad sobre rect\'angulo:} para $a_1 < a_2$, $b_1 < b_2$,
          \[
            P(a_1 < X \leq a_2,\, b_1 < Y \leq b_2)
            = F(a_2,b_2) - F(a_1,b_2) - F(a_2,b_1) + F(a_1,b_1).
          \]
  \end{enumerate}
\end{proposicion}

\vspace{4pt}
\begin{center}
  \begin{tikzpicture}[scale=0.85, >=Latex]
    \draw[-Latex, thick] (-0.3,0) -- (4.5,0) node[right,font=\small] {$X$};
    \draw[-Latex, thick] (0,-0.3) -- (0,3.5) node[above,font=\small] {$Y$};
    \fill[fill=defbg, draw=defborder, thick] (0,0) rectangle (3,2);
    \draw[dashed, thin, color=defborder] (3,0) -- (3,2) -- (0,2);
    \fill[color=defborder] (3,2) circle (2pt);
    \node[above right, font=\small\bfseries, color=defborder] at (3.05,2.05) {$(x_0, y_0)$};
    \node[below, font=\small] at (3,-0.05) {$x_0$};
    \node[left, font=\small] at (-0.05,2) {$y_0$};
    \node[font=\small\itshape, color=defborder] at (1.5,1.0) {$\{X\leq x_0,\,Y\leq y_0\}$};
    \node[font=\tiny, color=defborder] at (1.5,0.4) {prob.\ $=F_{X,Y}(x_0,y_0)$};
  \end{tikzpicture}
\end{center}
\begin{center}
  {\footnotesize\itshape La FDA conjunta $F_{X,Y}(x_0, y_0)$ asigna probabilidad al cuadrante
    inferior-izquierdo con esquina en $(x_0, y_0)$.}
\end{center}

\vspace{4pt}
\begin{center}
  \begin{tikzpicture}[scale=0.85, >=Latex]
    \draw[-Latex, thick] (-0.3,0) -- (4.7,0) node[right,font=\small] {$X$};
    \draw[-Latex, thick] (0,-0.3) -- (0,3.7) node[above,font=\small] {$Y$};
    \foreach \x/\y/\s in {1.0/1.0/3, 1.5/1.2/4, 2.0/1.5/5, 2.5/1.7/4, 3.0/2.0/3,
      1.2/1.8/3, 2.2/2.3/3, 0.8/1.6/2, 2.8/1.2/2, 1.7/0.8/2,
      3.3/2.5/2, 2.5/2.5/3, 1.5/2.5/2}
      \fill[fill=defborder, opacity=0.5] (\x,\y) circle (\s pt);
    \node[font=\small\itshape, color=defborder] at (2.0,3.2) {densidad conjunta $f_{X,Y}$};
    \draw[thick, color=noteborder] (0,-0.7) .. controls (1,-0.7) and (1.5,-1.5) .. (2.5,-1.6)
    .. controls (3.0,-1.5) and (3.8,-1.0) .. (4.5,-0.7);
    \node[font=\small\itshape, color=noteborder, right] at (4.5,-1.2) {$f_X(x)$ marginal};
    \draw[->, thin, gray, dashed] (1.5,1.2) -- (1.5,-0.6);
    \draw[->, thin, gray, dashed] (2.5,2.0) -- (2.5,-0.6);
    \node[font=\tiny\itshape, color=gray] at (3.7,-0.5) {integrar sobre $y$};
    \draw[thick, color=sectionbg] (-0.7,0) .. controls (-0.7,1) and (-1.5,1.5) .. (-1.6,2.0)
    .. controls (-1.5,2.5) and (-1.0,3.0) .. (-0.7,3.5);
    \node[font=\small\itshape, color=sectionbg, above, rotate=90] at (-1.4,2.0)
    {$f_Y(y)$ marginal};
    \draw[->, thin, gray, dashed] (1.5,1.2) -- (-0.6,1.2);
    \draw[->, thin, gray, dashed] (2.5,2.0) -- (-0.6,2.0);
  \end{tikzpicture}
\end{center}
\begin{center}
  {\footnotesize\itshape Las marginales $f_X$ y $f_Y$ son las \emph{proyecciones} de la
    densidad conjunta sobre cada eje.}
\end{center}

\begin{cuidado}
  \textbf{Las marginales no determinan la conjunta.} Conocer $F_X$ y $F_Y$ por separado no
  permite reconstruir $F_{X,Y}$. La conjunta agrega informaci\'on sobre la
  \emph{estructura de dependencia} entre las componentes. Dos vectores con las mismas
  marginales pueden tener conjuntas distintas.
\end{cuidado}

\subsecheader{5.4 \textperiodcentered\ Caso discreto}

Cuando $X, Y$ son v.a.\ discretas con valores en conjuntos numerables $I_X$ e $I_Y$,
la conjunta queda caracterizada por una \emph{funci\'on de masa}.

\begin{definicion}{Funci\'on de masa conjunta (PMF conjunta)}
  La \textbf{PMF conjunta} de $(X, Y)$ es:
  \[
    p_{X,Y}(x, y) = P(X = x,\, Y = y), \quad (x,y) \in I_X \times I_Y.
  \]
  Cumple: $p_{X,Y}(x,y) \geq 0$ \;y\; $\displaystyle\sum_{x\in I_X}\sum_{y\in I_Y} p_{X,Y}(x,y) = 1$.
\end{definicion}

\begin{definicion}{Funci\'on de masa marginal (PMF marginal)}
  Las \textbf{PMF marginales} se obtienen \emph{sumando} sobre la otra variable:
  \[
    p_X(x) = \sum_{y \in I_Y} p_{X,Y}(x, y), \qquad p_Y(y) = \sum_{x \in I_X} p_{X,Y}(x, y).
  \]
\end{definicion}

\begin{definicion}{FDA conjunta (caso discreto)}
  La \textbf{FDA conjunta} en el caso discreto se obtiene sumando la PMF en el cuadrante
  inferior-izquierdo:
  \[
    F_{X,Y}(x, y) = \sum_{x_i \leq x} \sum_{y_j \leq y} p_{X,Y}(x_i, y_j).
  \]
  Rec\'iprocamente, la PMF se recupera por diferencias finitas de la FDA.
\end{definicion}

\begin{ejemplo}{Dos dados: PMF conjunta y marginales}
  \textbf{Planteamiento.} Se lanzan dos dados justos de $6$ caras de manera independiente.
  Definimos:
  \begin{itemize}
    \item $X$ = resultado del primer dado, $X \in I_X = \{1,2,3,4,5,6\}$.
    \item $Y$ = resultado del segundo dado, $Y \in I_Y = \{1,2,3,4,5,6\}$.
  \end{itemize}
  El espacio muestral conjunto es $I_X \times I_Y$ con $|I_X \times I_Y| = 36$ pares posibles.
  \\[5pt]
  \textbf{(1) PMF conjunta.} Como los lanzamientos son \emph{independientes} y cada cara es
  equiprobable ($p_X(i) = 1/6$, $p_Y(j) = 1/6$), la PMF conjunta factoriza:
  \[
    p_{X,Y}(i,j) = P(X=i,\, Y=j) = P(X=i)\cdot P(Y=j) = \frac{1}{6}\cdot\frac{1}{6} = \frac{1}{36},
    \quad (i,j) \in I_X \times I_Y.
  \]
  \emph{Justificaci\'on:} la independencia permite descomponer la conjunta en producto de
  marginales. Si los dados estuvieran ``cargados'' o pegados (no independientes), no
  podr\'iamos hacer esto.
  \\[5pt]
  \textbf{(2) Verificaci\'on de normalizaci\'on.} La PMF debe sumar $1$:
  \[
    \sum_{i=1}^{6}\sum_{j=1}^{6} p_{X,Y}(i,j) = 36 \cdot \frac{1}{36} = 1. \;\checkmark
  \]
  \textbf{(3) Marginal de $X$.} ``Eliminamos'' $Y$ \emph{sumando} sobre todos sus valores:
  \[
    p_X(i) = \sum_{j=1}^{6} p_{X,Y}(i,j) = \underbrace{\frac{1}{36} + \frac{1}{36} + \cdots + \frac{1}{36}}_{6 \text{ veces}}
    = \frac{6}{36} = \frac{1}{6},
    \quad i \in \{1,\ldots,6\}.
  \]
  Por simetr\'ia, $p_Y(j) = 1/6$. \emph{Justificaci\'on:} sumar la fila $i$-\'esima recorre
  todas las opciones de $Y$ dado $X = i$, dejando solo la masa de probabilidad asociada a
  $X = i$ (sin importar qu\'e tom\'o $Y$).
  \\[5pt]
  \textbf{(4) Verificaci\'on de marginal v\'alida:} $\sum_{i=1}^6 p_X(i) = 6 \cdot 1/6 = 1$. \;\checkmark
  \\[5pt]
  \textbf{(5) Probabilidad sobre un evento.} Calculemos $P(X + Y = 7)$ (clave en juegos de dados):
  los pares $(i,j)$ con $i+j = 7$ son
  \[
    (1,6),\; (2,5),\; (3,4),\; (4,3),\; (5,2),\; (6,1) \quad \text{--- $6$ pares disjuntos.}
  \]
  Por aditividad de la PMF:
  \[
    P(X+Y = 7) = \sum_{(i,j):\, i+j=7} p_{X,Y}(i,j) = 6 \cdot \frac{1}{36} = \frac{1}{6} \approx 0.167.
  \]
  \textbf{(6) FDA conjunta evaluada.} Por ejemplo $F_{X,Y}(3, 4) = P(X\leq 3, Y\leq 4)$:
  son $3 \cdot 4 = 12$ pares, luego $F_{X,Y}(3,4) = 12/36 = 1/3$.
  \\[5pt]
  \textbf{(7) \textquestiondown Recuperan las marginales la conjunta?} Verificamos:
  $p_X(i)\,p_Y(j) = (1/6)(1/6) = 1/36 = p_{X,Y}(i,j)$. \;\checkmark
  \emph{En este caso s\'i,} porque sabemos que $X$ e $Y$ son independientes. En general, las
  marginales no determinan la conjunta.
  \\[5pt]
  \textbf{Lecci\'on clave.} Cuando $X, Y$ son independientes, la PMF conjunta es el producto
  de marginales y todo el c\'alculo se simplifica. Las marginales se obtienen sumando filas
  (para $X$) o columnas (para $Y$) en la cuadr\'icula de probabilidades.
\end{ejemplo}

\vspace{4pt}
\begin{center}
  \begin{tikzpicture}[scale=0.7]
    \foreach \i in {1,...,6}{
        \foreach \j in {1,...,6}{
            \fill[fill=defbg, draw=defborder, thin] (\i,\j) rectangle (\i+1,\j+1);
            \node[font=\tiny] at (\i+0.5,\j+0.5) {$\tfrac{1}{36}$};
          }
      }
    \foreach \i in {1,...,6}{
        \fill[fill=notebg, draw=noteborder, thin] (\i,0) rectangle (\i+1,0.7);
        \node[font=\tiny, color=noteborder] at (\i+0.5,0.35) {$\tfrac{1}{6}$};
        \node[font=\small\bfseries] at (\i+0.5,7.4) {$\i$};
      }
    \node[font=\small\itshape, color=noteborder, right] at (7.3,0.35) {$p_X(i)$};
    \foreach \j in {1,...,6}{
        \fill[fill=thmbg, draw=thmborder, thin] (0,\j) rectangle (0.7,\j+1);
        \node[font=\tiny, color=thmborder] at (0.35,\j+0.5) {$\tfrac{1}{6}$};
        \node[font=\small\bfseries] at (-0.4,\j+0.5) {$\j$};
      }
    \node[font=\small\itshape, color=thmborder] at (0.35,7.4) {$p_Y(j)$};
    \node[font=\small\bfseries] at (4,8) {$X$ (primer dado)};
    \node[font=\small\bfseries, rotate=90] at (-1.3,4) {$Y$ (segundo dado)};
  \end{tikzpicture}
\end{center}
\begin{center}
  {\footnotesize\itshape PMF conjunta del par $(X,Y)$ (cuadr\'icula central, todas $1/36$).
    Las marginales (barras laterales) se obtienen \emph{sumando} cada fila o columna.}
\end{center}

\begin{definicion}{PMF condicional (caso discreto)}
  Para $x \in I_X$ con $p_X(x) > 0$, la \textbf{PMF condicional} de $Y$ dado $X = x$ es:
  \[
    p_{Y \mid X}(y \mid x) = P(Y = y \mid X = x) = \frac{p_{X,Y}(x, y)}{p_X(x)},
    \quad y \in I_Y.
  \]
  Cumple: $p_{Y\mid X}(y\mid x) \geq 0$ y $\sum_{y\in I_Y} p_{Y\mid X}(y\mid x) = 1$
  (es una PMF v\'alida \emph{como funci\'on de $y$}, fijando $x$).
\end{definicion}

\begin{nota}
  \textbf{Intuici\'on.} La condicional $p_{Y\mid X}(\cdot \mid x)$ \emph{rebanada} la PMF
  conjunta a lo largo de la fila $X = x$ y la \emph{renormaliza} dividiendo por $p_X(x)$
  (la masa total de esa fila). Es la PMF de $Y$ \emph{dentro del subuniverso} donde $X = x$.
\end{nota}

\begin{ejemplo}{Suma de dos dados condicionada al primero}
  Sean $X$ = primer dado, $Y$ = segundo dado (independientes, justos), y $S = X + Y$ la suma.
  Calculemos la PMF condicional de $S$ dado $X$.
  \\[3pt]
  \textbf{(1) Conjunta de $(X, S)$:} para que $S = s$ dado $X = x$, debe ser $Y = s - x$, y eso
  ocurre si $s - x \in \{1, \ldots, 6\}$. Luego:
  \[
    p_{X,S}(x, s) = P(X=x,\, Y=s-x) = \begin{cases}
      \tfrac{1}{36} & x \in \{1,\ldots,6\},\; s-x \in \{1,\ldots,6\}, \\
      0             & \text{en otro caso}.
    \end{cases}
  \]
  \textbf{(2) PMF condicional de $S$ dado $X = x$:}
  \[
    p_{S\mid X}(s \mid x) = \frac{p_{X,S}(x, s)}{p_X(x)} = \frac{1/36}{1/6} = \frac{1}{6},
    \quad s \in \{x+1,\, x+2,\, \ldots,\, x+6\}.
  \]
  Es decir, dado $X = x$, la suma $S$ es uniforme en $\{x+1, \ldots, x+6\}$ (un desplazamiento
  de la distribuci\'on de $Y$).
  \\[3pt]
  \textbf{(3) Caso concreto:} si $X = 4$, entonces $S \in \{5, 6, 7, 8, 9, 10\}$ con probabilidad
  $1/6$ cada uno. Por ejemplo, $P(S = 7 \mid X = 4) = 1/6$.
  \\[3pt]
  \textbf{(4) Verificaci\'on:} $\sum_{s=x+1}^{x+6} p_{S\mid X}(s \mid x) = 6 \cdot 1/6 = 1$. \;\checkmark
\end{ejemplo}

\vspace{4pt}
\begin{center}
  \begin{tikzpicture}[scale=0.7]
    \foreach \i in {1,...,6}{
        \foreach \j in {1,...,6}{
            \ifnum\i=4
              \fill[fill=warnbg, draw=warnborder, thin] (\i,\j) rectangle (\i+1,\j+1);
              \node[font=\tiny, color=warnborder] at (\i+0.5,\j+0.5) {$\tfrac{1}{36}$};
            \else
              \fill[fill=defbg!50, draw=defborder!50, thin] (\i,\j) rectangle (\i+1,\j+1);
              \node[font=\tiny, color=gray] at (\i+0.5,\j+0.5) {$\tfrac{1}{36}$};
            \fi
          }
      }
    \foreach \i in {1,...,6}{
        \node[font=\small\bfseries] at (\i+0.5,7.4) {$\i$};
      }
    \foreach \j in {1,...,6}{
        \node[font=\small\bfseries] at (-0.4,\j+0.5) {$\j$};
      }
    \node[font=\small\bfseries] at (4,8) {$X$};
    \node[font=\small\bfseries, rotate=90] at (-1.3,4) {$Y$};
    \draw[thick, color=warnborder] (4,1) rectangle (5,7);
    \node[font=\small\bfseries, color=warnborder, right] at (7.3,4) {fila $X=4$:};
    \node[font=\tiny, color=warnborder, right] at (7.3,3.4) {masa total $1/6$};
    \draw[->, thick, color=warnborder] (5.2,4) -- (7.2,4);
    \node[font=\small\bfseries, color=warnborder, right] at (10.5,4) {$\Rightarrow$};
    \node[font=\small\bfseries, color=warnborder, right] at (11.2,4) {$p_{Y\mid X}(y\mid 4) = \tfrac{1}{6}$};
    \node[font=\tiny, color=warnborder, right] at (11.2,3.3) {(uniforme en $\{1,\ldots,6\}$)};
  \end{tikzpicture}
\end{center}
\begin{center}
  {\footnotesize\itshape Condicionar a $X = 4$ ``rebana'' la fila correspondiente y renormaliza
    su masa total ($1/6$) para obtener la PMF condicional de $Y$ (uniforme en $\{1,\ldots,6\}$,
    como era de esperarse por independencia).}
\end{center}

\subsecheader{5.5 \textperiodcentered\ Caso continuo}

Cuando $(X, Y)$ admite una funci\'on de densidad, la conjunta queda caracterizada por una
\emph{funci\'on de densidad}.

\begin{definicion}{Funci\'on de densidad conjunta (PDF conjunta)}
  $(X, Y)$ es \textbf{continuo} si existe $f_{X,Y}: \R^2 \to [0,\infty)$ tal que:
  \[
    F_{X,Y}(x, y) = \int_{-\infty}^{x}\!\int_{-\infty}^{y} f_{X,Y}(s, t)\,\mathrm{d}t\,\mathrm{d}s.
  \]
  Cumple: $f_{X,Y}(x,y) \geq 0$ \;y\; $\displaystyle\iint_{\R^2} f_{X,Y}(x,y)\,\mathrm{d}x\,\mathrm{d}y = 1$.\\
  \textbf{Probabilidad sobre un conjunto:} $P((X,Y) \in B) = \displaystyle\iint_B f_{X,Y}(x,y)\,\mathrm{d}x\,\mathrm{d}y$.
\end{definicion}

\begin{definicion}{Funci\'on de densidad marginal (PDF marginal)}
  Las \textbf{PDF marginales} se obtienen \emph{integrando} sobre la otra variable:
  \[
    f_X(x) = \int_{-\infty}^{\infty} f_{X,Y}(x, y)\,\mathrm{d}y, \qquad
    f_Y(y) = \int_{-\infty}^{\infty} f_{X,Y}(x, y)\,\mathrm{d}x.
  \]
\end{definicion}

\begin{definicion}{FDA conjunta (caso continuo)}
  La FDA se obtiene como integral doble de la PDF (definici\'on misma de continuidad).
  Rec\'iprocamente, donde $F_{X,Y}$ es diferenciable, la PDF se recupera por derivaci\'on:
  \[
    f_{X,Y}(x, y) = \frac{\partial^2 F_{X,Y}(x,y)}{\partial x\,\partial y}.
  \]
\end{definicion}

\begin{ejemplo}{Uniforme en el cuadrado unitario}
  \textbf{Planteamiento.} Sea $(X, Y)$ un punto elegido uniformemente al azar en el cuadrado
  $[0,1]^2$. ``Uniforme'' significa que ning\'un subregi\'on del cuadrado est\'a privilegiada:
  la probabilidad de que el punto caiga en una regi\'on es proporcional a su \emph{\'area}.
  \\[5pt]
  \textbf{(1) PDF conjunta.} La densidad uniforme sobre $[0,1]^2$ es constante:
  \[
    f_{X,Y}(x,y) = \begin{cases} 1 & 0 \leq x \leq 1,\; 0 \leq y \leq 1, \\ 0 & \text{en otro caso}.\end{cases}
  \]
  \emph{Justificaci\'on:} la PDF debe ser tal que $\iint f_{X,Y} = 1$; con soporte de \'area $1$
  (cuadrado unitario), la densidad constante que cumple esto es $f \equiv 1$ sobre el cuadrado.
  \\[5pt]
  \textbf{(2) Verificaci\'on de normalizaci\'on.}
  \[
    \iint_{\R^2} f_{X,Y}(x,y)\,\mathrm{d}x\,\mathrm{d}y
    = \iint_{[0,1]^2} 1\,\mathrm{d}x\,\mathrm{d}y
    = \int_0^1 \int_0^1 1\,\mathrm{d}y\,\mathrm{d}x
    = \int_0^1 1\,\mathrm{d}x = 1. \;\checkmark
  \]
  \textbf{(3) Marginal de $X$.} ``Eliminamos'' $Y$ integrando sobre todo su rango:
  \[
    f_X(x) = \int_{-\infty}^{\infty} f_{X,Y}(x,y)\,\mathrm{d}y = \int_0^1 1\,\mathrm{d}y = 1,
    \quad x \in [0,1].
  \]
  \emph{Justificaci\'on:} para cada $x$ fijo en $[0,1]$, $Y$ recorre $[0,1]$ uniformemente y
  contribuye total $1$. Fuera de $[0,1]$, la conjunta es $0$ y la marginal tambi\'en. \\
  Por simetr\'ia, $f_Y(y) = 1$ para $y \in [0,1]$. Por tanto:
  \[
    X \sim \mathrm{Unif}(0,1), \qquad Y \sim \mathrm{Unif}(0,1).
  \]
  \textbf{(4) FDA conjunta.} Para $(x,y) \in [0,1]^2$:
  \[
    F_{X,Y}(x,y) = \int_0^x \int_0^y 1\,\mathrm{d}t\,\mathrm{d}s = x\cdot y.
  \]
  Verificaci\'on: $f_{X,Y}(x,y) = \partial^2 F_{X,Y}/(\partial x\,\partial y) = 1$. \;\checkmark
  \\[5pt]
  \textbf{(5) Probabilidad sobre un conjunto: $P(X^2 + Y^2 \leq 1/4)$.}
  La regi\'on $B = \{(x,y): x^2+y^2 \leq 1/4\}$ intersectada con $[0,1]^2$ es el cuarto de
  c\'irculo de radio $1/2$ en el primer cuadrante. Como la densidad es $1$:
  \[
    P((X,Y) \in B) = \iint_B f_{X,Y}\,\mathrm{d}x\,\mathrm{d}y = \mathrm{\acute{a}rea}(B)
    = \frac{1}{4}\cdot \pi\!\left(\frac{1}{2}\right)^2 = \frac{\pi}{16} \approx 0.1963.
  \]
  \emph{Justificaci\'on:} con densidad uniforme $f \equiv 1$ sobre el cuadrado unitario, la
  probabilidad de cualquier subregi\'on \emph{es} su \'area. Esto es la traducci\'on continua
  de la ``regla cl\'asica'' $P(A) = |A|/|\Omega|$.
  \\[5pt]
  \textbf{(6) \textquestiondown Son $X$ e $Y$ independientes?} Verificamos si la conjunta
  factoriza:
  \[
    f_X(x)\,f_Y(y) = 1 \cdot 1 = 1 = f_{X,Y}(x,y) \quad \text{en } [0,1]^2. \;\checkmark
  \]
  S\'i factoriza, luego $X$ e $Y$ son \emph{independientes}.
  \\[5pt]
  \textbf{(7) Probabilidad de que $X > Y$.} Por simetr\'ia, $P(X > Y) = 1/2$. Verificaci\'on
  directa: el evento $\{X > Y\}$ es el tri\'angulo bajo la diagonal en $[0,1]^2$, de \'area $1/2$.
  \\[5pt]
  \textbf{Lecci\'on clave.} Para una distribuci\'on uniforme en una regi\'on $R \subseteq \R^2$
  con \'area finita, la PDF conjunta es $f = 1/\mathrm{\acute{a}rea}(R)$ sobre $R$, y la
  probabilidad de un sub-evento $B$ es $\mathrm{\acute{a}rea}(B \cap R)/\mathrm{\acute{a}rea}(R)$.
  Las marginales se obtienen integrando: para cada $x$, calcular la longitud del corte vertical
  del soporte.
\end{ejemplo}

\vspace{4pt}
\begin{center}
  \begin{tikzpicture}[scale=2.3, >=Latex]
    \draw[-Latex, thick] (-0.05,0) -- (1.3,0) node[right,font=\small] {$x$};
    \draw[-Latex, thick] (0,-0.05) -- (0,1.3) node[above,font=\small] {$y$};
    \fill[fill=defbg, draw=defborder, thick] (0,0) rectangle (1,1);
    \fill[fill=noteborder, opacity=0.55] (0,0) -- (0.5,0) arc (0:90:0.5) -- cycle;
    \draw[thick, color=noteborder] (0,0) -- (0.5,0) arc (0:90:0.5) -- cycle;
    \node[font=\small\bfseries, color=defborder] at (0.78,0.78) {$[0,1]^2$};
    \node[font=\small, color=noteborder] at (0.18,0.18) {$\tfrac{\pi}{16}$};
    \node[font=\tiny, color=noteborder] at (0.5,-0.12) {$x^2+y^2 \leq 1/4$};
    \node[below, font=\tiny] at (0.5,-0.02) {$\tfrac{1}{2}$};
    \node[below, font=\tiny] at (1,-0.02) {$1$};
    \node[left, font=\tiny] at (-0.02,0.5) {$\tfrac{1}{2}$};
    \node[left, font=\tiny] at (-0.02,1) {$1$};
  \end{tikzpicture}
\end{center}
\begin{center}
  {\footnotesize\itshape Soporte de $(X,Y) \sim \mathrm{Unif}([0,1]^2)$ (\'area $1$). Regi\'on
    sombreada \'ambar: $\{X^2+Y^2 \leq 1/4\}$, cuya probabilidad es su \'area: $\pi/16$.}
\end{center}

\begin{cuidado}
  \textbf{Independencia $\neq$ Descorrelaci\'on.} Si $X, Y$ son independientes, su covarianza es $0$
  y la PDF/PMF conjunta factoriza: $f_{X,Y} = f_X \cdot f_Y$. Pero el rec\'iproco \emph{no} es cierto
  en general: dos variables pueden tener covarianza $0$ \emph{sin} ser independientes (caso
  cl\'asico: $Y = X^2$ con $X \sim \mathrm{Unif}(-1, 1)$). La factorizaci\'on de la conjunta es la
  prueba estricta de independencia.
\end{cuidado}

\begin{definicion}{PDF condicional (caso continuo)}
  Para $x$ con $f_X(x) > 0$, la \textbf{PDF condicional} de $Y$ dado $X = x$ es:
  \[
    f_{Y \mid X}(y \mid x) = \frac{f_{X,Y}(x, y)}{f_X(x)},
    \quad y \in \R.
  \]
  Cumple: $f_{Y\mid X}(y\mid x) \geq 0$ y $\int_{\R} f_{Y\mid X}(y\mid x)\,\mathrm{d}y = 1$
  (es una PDF v\'alida \emph{como funci\'on de $y$}, fijando $x$).
\end{definicion}

\begin{nota}
  \textbf{Intuici\'on.} La condicional $f_{Y\mid X}(\cdot \mid x)$ es la \emph{rebanada vertical}
  de la densidad conjunta en $X = x$, renormalizada para integrar $1$. Es la densidad de $Y$
  cuando ya sabemos que $X = x$. \\
  \textbf{Cuidado con $P(X = x) = 0$:} en el caso continuo, condicionar a un punto exacto requiere
  esta definici\'on como \emph{l\'imite}, no como cociente directo de probabilidades.
\end{nota}

\begin{ejemplo}{Uniforme en un tri\'angulo}
  Sea $(X, Y)$ uniforme sobre el tri\'angulo $T = \{(x,y) : 0 \leq y \leq x \leq 1\}$.
  \\[3pt]
  \textbf{(1) PDF conjunta.} El \'area del tri\'angulo es $1/2$, luego:
  \[
    f_{X,Y}(x, y) = \begin{cases} 2 & 0 \leq y \leq x \leq 1, \\ 0 & \text{en otro caso.}\end{cases}
  \]
  \textbf{(2) Marginal de $X$.} Para $x \in [0,1]$, integrando $y$ sobre el rango v\'alido $[0, x]$:
  \[
    f_X(x) = \int_0^x 2\,\mathrm{d}y = 2x.
  \]
  Es una densidad triangular que favorece valores grandes de $X$.
  \\[3pt]
  \textbf{(3) PDF condicional de $Y$ dado $X = x$.} Para $x \in (0, 1]$:
  \[
    f_{Y\mid X}(y \mid x) = \frac{f_{X,Y}(x, y)}{f_X(x)} = \frac{2}{2x} = \frac{1}{x},
    \quad y \in [0, x].
  \]
  Es decir, $Y \mid X = x \sim \mathrm{Unif}(0, x)$: dado el primer valor $X = x$, $Y$ se
  distribuye uniformemente en $[0, x]$.
  \\[3pt]
  \textbf{(4) Verificaci\'on.}
  $\int_0^x \frac{1}{x}\,\mathrm{d}y = \frac{1}{x} \cdot x = 1$. \;\checkmark
  \\[3pt]
  \textbf{(5) Caso concreto.} Si $X = 0.7$, entonces $Y \sim \mathrm{Unif}(0, 0.7)$, con
  densidad $f_{Y\mid X}(y\mid 0.7) = 1/0.7 \approx 1.43$ sobre $[0, 0.7]$.
\end{ejemplo}

\vspace{4pt}
\begin{center}
  \begin{tikzpicture}[scale=2.3, >=Latex]
    \draw[-Latex, thick] (-0.05,0) -- (1.3,0) node[right,font=\small] {$x$};
    \draw[-Latex, thick] (0,-0.05) -- (0,1.3) node[above,font=\small] {$y$};
    \fill[fill=defbg, draw=defborder, thick] (0,0) -- (1,0) -- (1,1) -- cycle;
    \node[font=\small\bfseries, color=defborder] at (0.7,0.3) {$T:\, f_{X,Y}=2$};
    \draw[very thick, color=warnborder] (0.7,0) -- (0.7,0.7);
    \fill[color=warnborder] (0.7,0) circle (0.5pt);
    \fill[color=warnborder] (0.7,0.7) circle (0.5pt);
    \node[below, font=\tiny, color=warnborder] at (0.7,-0.05) {$X = 0.7$};
    \node[font=\tiny, color=warnborder, right] at (0.72,0.35) {rebanada};
    \draw[->, color=warnborder, thick] (0.85,0.35) -- (1.15,0.35);
    \node[font=\tiny, color=warnborder, right] at (1.16,0.35) {$Y\mid X = 0.7$};
    \node[font=\tiny, color=warnborder, right] at (1.16,0.25) {$\sim \mathrm{Unif}(0, 0.7)$};
    \node[below, font=\tiny] at (0,-0.02) {$0$};
    \node[below, font=\tiny] at (1,-0.02) {$1$};
  \end{tikzpicture}
\end{center}
\begin{center}
  {\footnotesize\itshape PDF condicional: la rebanada vertical en $X = 0.7$ tiene altura $2$
    y largo $0.7$, luego masa $1.4$. Renormalizando (dividiendo por $f_X(0.7) = 1.4$) se
    obtiene una distribuci\'on uniforme en $[0, 0.7]$.}
\end{center}

\subsecheader{5.6 \textperiodcentered\ Esperanza condicional}

La \textbf{esperanza condicional} $\E{Y \mid X = x}$ es la esperanza de $Y$ \emph{usando}
la distribuci\'on condicional $p_{Y\mid X}(\cdot \mid x)$ o $f_{Y\mid X}(\cdot \mid x)$.
Como funci\'on de $x$, define una variable aleatoria $\E{Y \mid X}$ que ``predice'' $Y$
a partir de $X$.

\begin{definicion}{Esperanza condicional (caso discreto)}
  Si $(X, Y)$ es discreto y $p_X(x) > 0$:
  \[
    \E{Y \mid X = x} = \sum_{y \in I_Y} y \cdot p_{Y \mid X}(y \mid x)
    = \frac{1}{p_X(x)} \sum_{y \in I_Y} y \cdot p_{X,Y}(x, y).
  \]
\end{definicion}

\begin{definicion}{Esperanza condicional (caso continuo)}
  Si $(X, Y)$ es continuo y $f_X(x) > 0$:
  \[
    \E{Y \mid X = x} = \int_{\R} y \cdot f_{Y \mid X}(y \mid x)\,\mathrm{d}y
    = \frac{1}{f_X(x)}\int_{\R} y \cdot f_{X,Y}(x, y)\,\mathrm{d}y.
  \]
\end{definicion}

\begin{nota}
  \textbf{Como variable aleatoria.} Definiendo $\E{Y \mid X}$ por la f\'ormula
  $\E{Y \mid X}(\omega) = \E{Y \mid X = X(\omega)}$, la esperanza condicional es ella misma
  una v.a.\ funci\'on de $X$. Su valor esperado es justamente $\E{Y}$ (ver propiedad de la torre).
\end{nota}

\begin{proposicion}{Propiedades de la esperanza condicional}
  Sean $X, Y$ v.a., $g: \R\to\R$ medible y $a, b \in \R$ constantes. Entonces:
  \begin{enumerate}
    \item \textbf{Linealidad:} $\E{aY_1 + bY_2 \mid X} = a\,\E{Y_1 \mid X} + b\,\E{Y_2 \mid X}$.
    \item \textbf{Constante:} $\E{c \mid X} = c$ para toda constante $c$.
    \item \textbf{Funci\'on de la condici\'on (sale fuera):} $\E{g(X)\, Y \mid X} = g(X)\, \E{Y \mid X}$.
          En particular $\E{g(X) \mid X} = g(X)$.
    \item \textbf{Independencia:} si $X \perp Y$, entonces $\E{Y \mid X} = \E{Y}$ (constante).
    \item \textbf{Propiedad de la torre (\emph{tower property}):}
          \[
            \E{\E{Y \mid X}} = \E{Y}.
          \]
          Esta es la \emph{ley de la esperanza total} en su forma m\'as general.
    \item \textbf{Monoton\'ia:} si $Y_1 \leq Y_2$ casi seguramente, $\E{Y_1 \mid X} \leq \E{Y_2 \mid X}$.
  \end{enumerate}
\end{proposicion}

\begin{ejemplo}{Esperanza condicional en el tri\'angulo}
  Continuamos con $(X, Y)$ uniforme en $T = \{0\leq y \leq x\leq 1\}$. Vimos que
  $Y \mid X = x \sim \mathrm{Unif}(0, x)$, luego:
  \[
    \E{Y \mid X = x} = \frac{0 + x}{2} = \frac{x}{2}, \quad x \in [0, 1].
  \]
  Por tanto, como variable aleatoria, $\E{Y \mid X} = X/2$.
  \\[3pt]
  \textbf{Verificaci\'on por la torre.} Calculamos $\E{Y}$ de dos maneras:
  \begin{itemize}
    \item \textbf{Por la torre:} $\E{Y} = \E{\E{Y\mid X}} = \E{X/2}$.
          Recordando que $f_X(x) = 2x$, $\E{X} = \int_0^1 x\cdot 2x\,\mathrm{d}x = 2/3$.
          Luego $\E{Y} = (1/2) \cdot (2/3) = 1/3$.
    \item \textbf{Directo:} $\E{Y} = \int_0^1 y \cdot f_Y(y)\,\mathrm{d}y$ con $f_Y(y) = 2(1-y)$:
          $\E{Y} = \int_0^1 y\cdot 2(1-y)\,\mathrm{d}y = \int_0^1 (2y - 2y^2)\,\mathrm{d}y = 1 - 2/3 = 1/3$. \;\checkmark
  \end{itemize}
  Ambos m\'etodos coinciden.
\end{ejemplo}

\begin{nota}
  \textbf{Resumen:} Conjunta $\to$ Marginal: suma/integraci\'on. PMF/PDF $\leftrightarrow$ FDA:
  suma/integraci\'on o derivaci\'on. Marginal $\not\Rightarrow$ Conjunta: falta info de dependencia.
  \\[3pt]
  \textbf{Condicionales:} rebanan la conjunta y renormalizan. $\E{Y\mid X}$ es v.a.\ funci\'on
  de $X$; la torre ($\E{\E{Y\mid X}} = \E{Y}$) permite calcular $\E{Y}$ ``por capas''.
\end{nota}

% =============================================================================
\secheader{6}{Proceso de Poisson}

\subsecheader{6.1 \textperiodcentered\ Procesos estoc\'asticos y de conteo}

\begin{definicion}{Proceso estoc\'astico}
  Un \textbf{proceso estoc\'astico} es una familia de variables aleatorias
  $\{X(t),\, t\in T\}$ definidas sobre el mismo espacio de probabilidad, con $T$ conjunto de
  \'indices (t\'ipicamente tiempo).
  \begin{itemize}
    \item Si $T = \{0, 1, 2, \ldots\}$: tiempo discreto.
    \item Si $T = [0, \infty)$: tiempo continuo.
  \end{itemize}
\end{definicion}

\begin{definicion}{Proceso de conteo}
  $\{N(t),\, t\geq 0\}$ es un \textbf{proceso de conteo} si:
  \begin{enumerate}
    \item $N(t) \in \Nset_0$.
    \item $N(t)$ es no decreciente.
    \item Para $s < t$, $N(t) - N(s) = $ \# eventos en $(s, t]$.
  \end{enumerate}
  T\'ipicamente $N(0) = 0$.
\end{definicion}

\begin{center}
  \begin{tikzpicture}[scale=0.95, >=Latex]
    \draw[-Latex, thick] (-0.2,0) -- (8.5,0) node[right,font=\small] {$t$};
    \draw[-Latex, thick] (0,-0.2) -- (0,4.2) node[above,font=\small] {$N(t)$};
    \draw[very thick, color=sectionbg]
    (0,0) -- (1.2,0) -- (1.2,1) -- (2.5,1) -- (2.5,2) --
    (3.3,2) -- (3.3,3) -- (5.0,3) -- (5.0,4) -- (8.2,4);
    \foreach \x/\y in {1.2/1, 2.5/2, 3.3/3, 5.0/4} {
        \fill[color=sectionbg] (\x,\y) circle (2.5pt);
        \draw[dashed, color=gray!50, thin] (\x,0) -- (\x,\y);
      }
    \foreach \y in {1,2,3,4} {
        \node[font=\tiny, left, color=gray] at (0,\y) {$\y$};
      }
  \end{tikzpicture}
\end{center}

\subsecheader{6.2 \textperiodcentered\ Incrementos independientes y estacionarios}

Un \textbf{incremento} $N(t) - N(s)$ del proceso es el n\'umero de eventos ocurridos en el
intervalo $(s, t]$. Las dos propiedades siguientes describen \emph{c\'omo se relacionan}
los incrementos sobre distintos intervalos.

\begin{definicion}{Incrementos independientes}
  $\{N(t)\}$ tiene \textbf{incrementos independientes} si para cualquier partici\'on
  $0 \leq t_0 < t_1 < t_2 < \cdots < t_k$, las variables aleatorias
  \[
    N(t_1) - N(t_0),\quad N(t_2) - N(t_1),\quad \ldots,\quad N(t_k) - N(t_{k-1})
  \]
  son \emph{mutuamente independientes}.
  \\[3pt]
  \emph{Interpretaci\'on:} lo que ocurre en un intervalo \emph{no aporta informaci\'on}
  sobre lo que ocurre en cualquier otro intervalo \emph{disjunto}. Cada intervalo es un
  ``mundo aparte'' estad\'isticamente.
\end{definicion}

\vspace{4pt}
\begin{center}
  \begin{tikzpicture}[scale=0.9, >=Latex]
    \draw[-Latex, thick] (-0.2,0) -- (10,0) node[right,font=\small] {$t$};
    \fill[fill=defbg, opacity=0.7] (1,-0.18) rectangle (3.5,0.18);
    \fill[fill=notebg, opacity=0.7] (5,-0.18) rectangle (8,0.18);
    \draw[thick, color=defborder] (1,-0.2) -- (1,0.2);
    \draw[thick, color=defborder] (3.5,-0.2) -- (3.5,0.2);
    \draw[thick, color=noteborder] (5,-0.2) -- (5,0.2);
    \draw[thick, color=noteborder] (8,-0.2) -- (8,0.2);
    \node[below, font=\small, color=defborder] at (1,-0.25) {$s$};
    \node[below, font=\small, color=defborder] at (3.5,-0.25) {$t$};
    \node[below, font=\small, color=noteborder] at (5,-0.25) {$u$};
    \node[below, font=\small, color=noteborder] at (8,-0.25) {$v$};
    \foreach \x in {1.5, 2.5, 3.0, 5.5, 6.8, 7.5} {
        \fill[color=sectionbg] (\x,0) circle (2.5pt);
      }
    \node[above, font=\small\bfseries, color=defborder] at (2.25,0.55)
    {$N(t) - N(s) = 3$};
    \node[above, font=\small\bfseries, color=noteborder] at (6.5,0.55)
    {$N(v) - N(u) = 3$};
    \draw[<->, thick, color=gray, dashed] (2.25,1.05) to[bend left=20] (6.5,1.05);
    \node[font=\small\itshape, color=gray, above] at (4.4,1.45) {\textbf{independientes}};
    \node[font=\tiny\itshape, color=gray] at (4.4,-0.8) {intervalos disjuntos};
  \end{tikzpicture}
\end{center}
\begin{center}
  {\footnotesize\itshape Incrementos independientes: el conteo en $(s, t]$ no informa sobre
    el conteo en $(u, v]$ (intervalos disjuntos). Cada intervalo es ``un mundo aparte''.}
\end{center}

\begin{definicion}{Incrementos estacionarios}
  $\{N(t)\}$ tiene \textbf{incrementos estacionarios} si la distribuci\'on de
  $N(t+s) - N(t)$ depende solo de la longitud $s$ del intervalo, no de su posici\'on $t$:
  \[
    N(t + s) - N(t) \;\stackrel{d}{=}\; N(s) - N(0), \qquad \forall\, t, s \geq 0.
  \]
  \emph{Interpretaci\'on:} la ``tasa'' del proceso es \emph{constante en el tiempo}; intervalos
  de la misma longitud tienen la misma distribuci\'on de conteos sin importar d\'onde
  est\'en colocados.
\end{definicion}

\vspace{4pt}
\begin{center}
  \begin{tikzpicture}[scale=0.9, >=Latex]
    \draw[-Latex, thick] (-0.2,0) -- (10,0) node[right,font=\small] {$t$};
    \fill[fill=defbg, opacity=0.7] (0.5,-0.18) rectangle (2,0.18);
    \fill[fill=defbg, opacity=0.7] (3.5,-0.18) rectangle (5,0.18);
    \fill[fill=defbg, opacity=0.7] (6.5,-0.18) rectangle (8,0.18);
    \foreach \xa/\xb in {0.5/2, 3.5/5, 6.5/8} {
        \draw[thick, color=defborder] (\xa,-0.2) -- (\xa,0.2);
        \draw[thick, color=defborder] (\xb,-0.2) -- (\xb,0.2);
      }
    \draw[<->, color=defborder, thin] (0.5,-0.55) -- (2,-0.55);
    \draw[<->, color=defborder, thin] (3.5,-0.55) -- (5,-0.55);
    \draw[<->, color=defborder, thin] (6.5,-0.55) -- (8,-0.55);
    \node[below, font=\tiny, color=defborder] at (1.25,-0.65) {longitud $s$};
    \node[below, font=\tiny, color=defborder] at (4.25,-0.65) {longitud $s$};
    \node[below, font=\tiny, color=defborder] at (7.25,-0.65) {longitud $s$};
    \foreach \x in {0.7, 1.4, 1.8, 3.0, 3.8, 4.6, 6.0, 6.7, 7.6, 9.0} {
        \fill[color=sectionbg] (\x,0) circle (2pt);
      }
    \node[above, font=\tiny\bfseries, color=defborder] at (1.25,0.3) {$N_1$};
    \node[above, font=\tiny\bfseries, color=defborder] at (4.25,0.3) {$N_2$};
    \node[above, font=\tiny\bfseries, color=defborder] at (7.25,0.3) {$N_3$};
    \node[above, font=\small\itshape, color=defborder] at (4.25,0.7)
    {$N_1 \stackrel{d}{=} N_2 \stackrel{d}{=} N_3 \sim \mathrm{Poisson}(\lambda s)$};
  \end{tikzpicture}
\end{center}
\begin{center}
  {\footnotesize\itshape Incrementos estacionarios: tres intervalos de la \emph{misma longitud}
    $s$ ubicados en posiciones distintas; sus conteos tienen \emph{id\'entica distribuci\'on}.
    En un proceso de Poisson de tasa $\lambda$, todos $\sim \mathrm{Poisson}(\lambda s)$.}
\end{center}

\begin{nota}
  \textbf{Diferencia clave.} Las dos propiedades describen cosas distintas:
  \begin{itemize}
    \item \textbf{Independientes:} dos intervalos \emph{disjuntos cualesquiera} no se influyen.
    \item \textbf{Estacionarios:} el mismo intervalo \emph{trasladado} mantiene su distribuci\'on.
  \end{itemize}
  Un proceso puede tener una sin la otra. El proceso de Poisson \emph{homog\'eneo} las cumple
  ambas; el proceso de Poisson \emph{no-homog\'eneo} (con tasa $\lambda(t)$ variable) tiene
  incrementos independientes pero \emph{no} estacionarios.
\end{nota}

\vspace{6pt}
\textbf{Tres escenarios al trabajar con dos intervalos.}
Para un proceso de Poisson de tasa $\lambda = 2$/min, fijamos dos intervalos y analizamos
sus incrementos seg\'un la posici\'on relativa. La regla general:
\emph{descomponer en piezas disjuntas, aplicar incrementos independientes, sumar Poissons}.

\begin{ejemplo}{Caso 1: intervalos disjuntos}
  \textbf{Intervalos:} $(0, 3]$ y $(5, 8]$ minutos. $\lambda = 2$/min.
  \\[3pt]
  \begin{center}
    \begin{tikzpicture}[scale=0.75, >=Latex]
      \draw[-Latex, thick] (-0.2,0) -- (10,0) node[right,font=\small] {$t$};
      \fill[fill=defbg, opacity=0.7] (0,-0.18) rectangle (3,0.18);
      \fill[fill=notebg, opacity=0.7] (5,-0.18) rectangle (8,0.18);
      \draw[thick, color=defborder] (0,-0.2) -- (0,0.2);
      \draw[thick, color=defborder] (3,-0.2) -- (3,0.2);
      \draw[thick, color=noteborder] (5,-0.2) -- (5,0.2);
      \draw[thick, color=noteborder] (8,-0.2) -- (8,0.2);
      \node[below, font=\tiny, color=defborder] at (0,-0.25) {$0$};
      \node[below, font=\tiny, color=defborder] at (3,-0.25) {$3$};
      \node[below, font=\tiny, color=noteborder] at (5,-0.25) {$5$};
      \node[below, font=\tiny, color=noteborder] at (8,-0.25) {$8$};
      \node[above, font=\tiny\bfseries, color=defborder] at (1.5,0.25) {$N_A$};
      \node[above, font=\tiny\bfseries, color=noteborder] at (6.5,0.25) {$N_B$};
    \end{tikzpicture}
  \end{center}
  \textbf{An\'alisis.} Como $(0,3] \cap (5,8] = \emptyset$, los incrementos son
  \emph{independientes}:
  \[
    N_A := N(3) - N(0) \sim \mathrm{Poisson}(\lambda \cdot 3) = \mathrm{Poisson}(6),
    \quad N_B := N(8) - N(5) \sim \mathrm{Poisson}(6), \quad N_A \perp N_B.
  \]
  \textbf{Conjunta.} La probabilidad de un evento conjunto factoriza:
  \[
    P(N_A = 2,\, N_B = 4) = P(N_A = 2) \cdot P(N_B = 4)
    = \frac{6^2 e^{-6}}{2!}\cdot\frac{6^4 e^{-6}}{4!} \approx 0.0446 \cdot 0.1339 \approx 0.0060.
  \]
\end{ejemplo}

\begin{ejemplo}{Caso 2: intervalos que se intersectan}
  \textbf{Intervalos:} $(0, 5]$ y $(3, 8]$ minutos (se solapan en $(3, 5]$). $\lambda = 2$/min.
  \\[3pt]
  \begin{center}
    \begin{tikzpicture}[scale=0.75, >=Latex]
      \draw[-Latex, thick] (-0.2,0) -- (10,0) node[right,font=\small] {$t$};
      \fill[fill=defbg, opacity=0.6] (0,0.05) rectangle (5,0.35);
      \fill[fill=notebg, opacity=0.6] (3,-0.35) rectangle (8,-0.05);
      \fill[fill=warnbg, opacity=0.7] (3,-0.05) rectangle (5,0.05);
      \draw[thick, color=defborder] (0,0.05) -- (0,0.35);
      \draw[thick, color=defborder] (5,0.05) -- (5,0.35);
      \draw[thick, color=noteborder] (3,-0.35) -- (3,-0.05);
      \draw[thick, color=noteborder] (8,-0.35) -- (8,-0.05);
      \node[above, font=\tiny\bfseries, color=defborder] at (2.5,0.4) {$(0, 5]$};
      \node[below, font=\tiny\bfseries, color=noteborder] at (5.5,-0.4) {$(3, 8]$};
      \node[font=\tiny\bfseries, color=warnborder] at (4,-0.6) {$\uparrow$ overlap $(3,5]$};
      \node[below, font=\tiny] at (0,-0.7) {$0$};
      \node[below, font=\tiny] at (3,-0.7) {$3$};
      \node[below, font=\tiny] at (5,-0.7) {$5$};
      \node[below, font=\tiny] at (8,-0.7) {$8$};
    \end{tikzpicture}
  \end{center}
  \textbf{Estrategia: descomponer en tres piezas disjuntas.}
  \[
    A := N(3) - N(0),\quad B := N(5) - N(3),\quad C := N(8) - N(5).
  \]
  Por incrementos independientes y estacionarios:
  \[
    A \sim \mathrm{Poisson}(6),\quad B \sim \mathrm{Poisson}(4),\quad C \sim \mathrm{Poisson}(6),
    \quad A, B, C \text{ mutuamente independientes}.
  \]
  Luego:
  \begin{itemize}
    \item $N(5) - N(0) = A + B \sim \mathrm{Poisson}(10)$ \;(suma de Poissons indep.).
    \item $N(8) - N(3) = B + C \sim \mathrm{Poisson}(10)$.
    \item \textbf{NO son independientes}: comparten $B$. La covarianza es:
          \[
            \Cov{A+B}{B+C} = \Var{B} = 4.
          \]
  \end{itemize}
  \textbf{Probabilidad conjunta.} Para calcular $P(N(5) - N(0) = k_1,\; N(8) - N(3) = k_2)$
  hay que sumar sobre los valores compartidos de $B$:
  \[
    P(A+B = k_1,\, B+C = k_2)
    = \sum_{b=0}^{\min(k_1, k_2)} P(B = b)\, P(A = k_1 - b)\, P(C = k_2 - b).
  \]
\end{ejemplo}

\begin{ejemplo}{Caso 3: un intervalo dentro del otro}
  \textbf{Intervalos:} $(0, 8]$ contiene a $(2, 5]$. $\lambda = 2$/min.
  \\[3pt]
  \begin{center}
    \begin{tikzpicture}[scale=0.75, >=Latex]
      \draw[-Latex, thick] (-0.2,0) -- (10,0) node[right,font=\small] {$t$};
      \fill[fill=defbg, opacity=0.5] (0,-0.4) rectangle (8,0.4);
      \fill[fill=warnbg, opacity=0.8] (2,-0.18) rectangle (5,0.18);
      \draw[thick, color=defborder] (0,-0.4) -- (0,0.4);
      \draw[thick, color=defborder] (8,-0.4) -- (8,0.4);
      \draw[thick, color=warnborder] (2,-0.2) -- (2,0.2);
      \draw[thick, color=warnborder] (5,-0.2) -- (5,0.2);
      \node[below, font=\tiny] at (0,-0.45) {$0$};
      \node[below, font=\tiny] at (2,-0.45) {$2$};
      \node[below, font=\tiny] at (5,-0.45) {$5$};
      \node[below, font=\tiny] at (8,-0.45) {$8$};
      \node[above, font=\tiny\bfseries, color=defborder] at (4,0.5) {exterior $(0,8]$};
      \node[font=\tiny\bfseries, color=warnborder] at (3.5,0.0) {interior $(2,5]$};
    \end{tikzpicture}
  \end{center}
  \textbf{Estrategia: descomponer en tres piezas disjuntas.}
  \[
    A_1 := N(2) - N(0),\quad A_2 := N(5) - N(2),\quad A_3 := N(8) - N(5).
  \]
  Por incrementos independientes y estacionarios:
  \[
    A_1 \sim \mathrm{Poisson}(4),\quad A_2 \sim \mathrm{Poisson}(6),\quad A_3 \sim \mathrm{Poisson}(6),
    \quad \text{mutuamente independientes}.
  \]
  Luego:
  \begin{itemize}
    \item Conteo del intervalo grande: $N(8) - N(0) = A_1 + A_2 + A_3 \sim \mathrm{Poisson}(16)$.
    \item Conteo del intervalo peque\~no: $N(5) - N(2) = A_2 \sim \mathrm{Poisson}(6)$.
    \item \textbf{NO son independientes}: el grande \emph{contiene} al peque\~no, luego
          $N(8) - N(0) \geq N(5) - N(2)$ siempre. La covarianza:
          $\Cov{A_1+A_2+A_3}{A_2} = \Var{A_2} = 6$.
    \item Pero la \emph{diferencia} grande $-$ peque\~no $= A_1 + A_3 \sim \mathrm{Poisson}(10)$ \emph{s\'i} es
          independiente de $A_2 = $ peque\~no. Las piezas \emph{fuera} del interior
          son independientes del interior.
  \end{itemize}
\end{ejemplo}

\begin{nota}
  \textbf{Receta universal.} Frente a varios intervalos relacionados:
  \begin{enumerate}
    \item Identifica las \emph{piezas at\'omicas disjuntas} (los intervalos m\'as peque\~nos
          que generan todo lo dem\'as por uniones).
    \item Por estacionariedad, cada pieza es Poisson de tasa $\lambda \cdot$(longitud).
    \item Por independencia de incrementos, las piezas at\'omicas son mutuamente independientes.
    \item Cualquier conteo de inter\'es se expresa como suma de piezas; las sumas de Poissons
          independientes son Poissons.
    \item La dependencia entre dos conteos compuestos se mide por las piezas \emph{compartidas}.
  \end{enumerate}
\end{nota}

\subsecheader{6.3 \textperiodcentered\ Proceso de Poisson homog\'eneo}

\begin{definicion}{Proceso de Poisson homog\'eneo}
  $\{N(t)\}$ es un \textbf{proceso de Poisson} con tasa $\lambda > 0$ si:
  \begin{enumerate}
    \item $N(0) = 0$.
    \item Tiene incrementos independientes.
    \item $N(t+s) - N(t) \sim \mathrm{Poisson}(\lambda s)$.
  \end{enumerate}
\end{definicion}

\begin{proposicion}{Distribuci\'on de $N(t)$}
  En un proceso de Poisson con tasa $\lambda > 0$, para todo $t \geq 0$:
  \[
    N(t) \;\sim\; \mathrm{Poisson}(\lambda t).
  \]
  En particular:
  \[
    P(N(t) = k) = \frac{(\lambda t)^k e^{-\lambda t}}{k!}, \quad k \in \Nset_0,
    \qquad \E{N(t)} = \Var{N(t)} = \lambda t.
  \]
\end{proposicion}

\begin{demostracion}
  De la propiedad (1), $N(0) = 0$. Aplicando la propiedad (3) con $s := t$ y $t := 0$:
  \[
    N(t) - N(0) \;\sim\; \mathrm{Poisson}(\lambda t),
  \]
  pero $N(t) - N(0) = N(t) - 0 = N(t)$, de donde $N(t) \sim \mathrm{Poisson}(\lambda t)$.
  \\[3pt]
  Los momentos $\E{N(t)} = \Var{N(t)} = \lambda t$ se siguen de las propiedades de la
  distribuci\'on Poisson (equidispersi\'on, Secci\'on 4.5).
\end{demostracion}

\begin{nota}
  \textbf{Interpretaci\'on:} en un intervalo de longitud $t$, el n\'umero de eventos sigue
  Poisson con par\'ametro $\lambda t$. Por tanto en promedio ocurren $\lambda t$ eventos en
  un intervalo de longitud $t$, lo que justifica llamar a $\lambda$ la \emph{tasa}
  (eventos por unidad de tiempo).
\end{nota}

\subsecheader{6.4 \textperiodcentered\ Tiempos entre eventos}

\begin{proposicion}{$T_i \sim$ Exponencial}
  En un proceso de Poisson de tasa $\lambda$, los tiempos entre eventos $T_1, T_2, \ldots$
  son i.i.d.\ $\sim \mathrm{Exp}(\lambda)$.
\end{proposicion}

\begin{demostracion}
  $P(T_1 > t) = P(N(t)=0) = e^{-\lambda t}$, luego $T_1 \sim \mathrm{Exp}(\lambda)$.
  Por incrementos independientes y estacionarios, $T_2$ es independiente de $T_1$ y tambi\'en
  $\mathrm{Exp}(\lambda)$. Inducci\'on completa.
\end{demostracion}

\begin{center}
  \begin{tikzpicture}[scale=0.9, >=Latex]
    \draw[-Latex, thick] (-0.2,0) -- (9,0) node[right,font=\small] {$t$};
    \foreach \x/\n in {1.2/1, 2.5/2, 3.3/3, 5.0/4, 7.5/5} {
        \fill[color=sectionbg] (\x,0) circle (2.5pt);
        \node[above, font=\tiny\bfseries, color=sectionbg] at (\x,0.12) {$S_\n$};
      }
    \draw[{Latex[length=2.5pt]}-{Latex[length=2.5pt]}, color=defborder, thick]
    (0,-0.55) -- (1.2,-0.55) node[midway,below,font=\tiny\itshape,color=defborder] {$T_1$};
    \draw[{Latex[length=2.5pt]}-{Latex[length=2.5pt]}, color=defborder, thick]
    (1.2,-0.55) -- (2.5,-0.55) node[midway,below,font=\tiny\itshape,color=defborder] {$T_2$};
    \draw[{Latex[length=2.5pt]}-{Latex[length=2.5pt]}, color=defborder, thick]
    (2.5,-0.55) -- (3.3,-0.55) node[midway,below,font=\tiny\itshape,color=defborder] {$T_3$};
    \draw[{Latex[length=2.5pt]}-{Latex[length=2.5pt]}, color=defborder, thick]
    (3.3,-0.55) -- (5.0,-0.55) node[midway,below,font=\tiny\itshape,color=defborder] {$T_4$};
    \draw[{Latex[length=2.5pt]}-{Latex[length=2.5pt]}, color=defborder, thick]
    (5.0,-0.55) -- (7.5,-0.55) node[midway,below,font=\tiny\itshape,color=defborder] {$T_5$};
    \node[font=\tiny\itshape, color=defborder] at (4.5,-1.0)
    {$T_k \sim \mathrm{Exp}(\lambda)$ i.i.d.\,;\quad $S_n = T_1+\cdots+T_n$};
  \end{tikzpicture}
\end{center}
\begin{center}
  {\footnotesize\itshape Tiempos entre llegadas $T_k$ y tiempos de las $n$-\'esimas llegadas $S_n$.}
\end{center}

\begin{cuidado}
  \textbf{P\'erdida de memoria es contraintuitivo.} Si llevas $10$ min esperando un evento
  Exp$(\lambda)$ sin que ocurra, la probabilidad de esperar otros $5$ min es \emph{la misma}
  que esperar $5$ min desde cero. NO te ``acercas'' al evento por haber esperado.
  Este es el modelo de \emph{fallo aleatorio puro} (sin desgaste). Si el sistema tiene
  envejecimiento real (componentes que se deterioran, baterías que se gastan), la Exponencial
  NO es apropiada --- se necesita una distribuci\'on con tasa de fallo creciente como Weibull.
\end{cuidado}

\subsecheader{6.5 \textperiodcentered\ Competencia de tasa}

\begin{proposicion}{Competencia entre dos procesos de Poisson}
  Sean $N_1, N_2$ procesos independientes con tasas $\lambda_1, \lambda_2$. Para los primeros
  eventos $T^{(i)}$:
  \begin{enumerate}
    \item $T = \min(T^{(1)}, T^{(2)}) \sim \mathrm{Exp}(\lambda_1 + \lambda_2)$.
    \item $P(T^{(1)} < T^{(2)}) = \lambda_1/(\lambda_1 + \lambda_2)$.
  \end{enumerate}
\end{proposicion}

\begin{demostracion}
  $P(T > t) = e^{-\lambda_1 t}\,e^{-\lambda_2 t} = e^{-(\lambda_1+\lambda_2)t}$.
  Por integraci\'on: $P(T^{(1)} < T^{(2)}) = \int_0^\infty \lambda_1 e^{-\lambda_1 t} e^{-\lambda_2 t}\,\mathrm{d}t
  = \lambda_1/(\lambda_1+\lambda_2)$.
\end{demostracion}

\begin{ejemplo}{Dos cajas en un banco}
  Llegan clientes de dos tipos a un banco, de manera independiente:
  \begin{itemize}
    \item Tipo $A$ (caja preferencial) a tasa $\lambda_A = 6$/hora.
    \item Tipo $B$ (caja regular) a tasa $\lambda_B = 4$/hora.
  \end{itemize}
  \textbf{(1) Tiempo hasta el pr\'oximo cliente (cualquier tipo):}
  $T = \min(T^{(A)}, T^{(B)}) \sim \mathrm{Exp}(6+4) = \mathrm{Exp}(10)$. Por tanto:
  \[
    \E{T} = \frac{1}{\lambda_A + \lambda_B} = \frac{1}{10}\text{ h} = 6 \text{ min}.
  \]
  \textbf{(2) Probabilidad de que el pr\'oximo sea de tipo $A$:}
  \[
    P(T^{(A)} < T^{(B)}) = \frac{\lambda_A}{\lambda_A + \lambda_B} = \frac{6}{10} = 0.6.
  \]
  \textbf{(3) Probabilidad de esperar m\'as de $5$ minutos:}
  \[
    P(T > 5/60) = e^{-10\cdot 5/60} = e^{-5/6} \approx 0.435.
  \]
  \textbf{(4) Interpretaci\'on:} la mezcla de los dos procesos act\'ua como un solo proceso de
  Poisson con tasa suma. Cada cliente que llega es independientemente tipo $A$ con probabilidad
  $\lambda_A/(\lambda_A+\lambda_B) = 0.6$ o tipo $B$ con $0.4$, sin importar las llegadas previas.
\end{ejemplo}

\begin{center}
  \begin{tikzpicture}[scale=0.95, >=Latex]
    \draw[-Latex, thick] (-0.2,1.5) -- (8.5,1.5) node[right,font=\small] {tipo $1$};
    \foreach \x in {1.6, 4.0, 6.7} {
        \fill[color=defborder] (\x,1.5) circle (3pt);
      }
    \node[font=\small\bfseries, color=defborder, left] at (-0.3,1.5) {$\lambda_1$};
    \draw[-Latex, thick] (-0.2,0) -- (8.5,0) node[right,font=\small] {tipo $2$};
    \foreach \x in {0.9, 2.8, 5.2, 7.5} {
        \fill[color=noteborder] (\x,0) circle (3pt);
      }
    \node[font=\small\bfseries, color=noteborder, left] at (-0.3,0) {$\lambda_2$};
    \draw[->, very thick, color=sectionbg] (0.9, -0.6) -- (0.9, -0.2);
    \node[font=\small\bfseries, color=sectionbg, below] at (0.9,-0.7)
    {primer evento (tipo $2$)};
  \end{tikzpicture}
\end{center}

\subsecheader{6.6 \textperiodcentered\ Tiempo de la $n$-\'esima llegada}

\begin{proposicion}{$S_n \sim$ Erlang}
  $S_n = T_1 + \cdots + T_n \sim \mathrm{Erlang}(n, \lambda)$:
  \[
    f_{S_n}(t) = \frac{\lambda^n t^{n-1} e^{-\lambda t}}{(n-1)!}, \quad t \geq 0.
  \]
  $\E{S_n} = n/\lambda$, $\Var{S_n} = n/\lambda^2$.
\end{proposicion}

\begin{demostracion}
  Conexi\'on clave: $\{S_n \leq t\} = \{N(t) \geq n\}$ (la $n$-\'esima llegada ocurre antes
  de $t$ si y solo si han ocurrido al menos $n$ eventos hasta $t$). Entonces:
  \[
    F_{S_n}(t) = P(N(t)\geq n) = \sum_{k=n}^\infty \frac{(\lambda t)^k e^{-\lambda t}}{k!}.
  \]
  Derivando respecto a $t$ y simplificando (los t\'erminos se cancelan en cascada):
  \[
    f_{S_n}(t) = \frac{\lambda^n t^{n-1} e^{-\lambda t}}{(n-1)!}.
  \]
\end{demostracion}

\begin{ejemplo}{Tiempo hasta la quinta llamada}
  En un call center con tasa $\lambda = 5$ llamadas/min, el tiempo $S_5$ hasta la quinta llamada
  sigue una Erlang($5, 5$).
  \\[3pt]
  \textbf{Tiempo medio:} $\E{S_5} = 5/5 = 1$ min.\\
  \textbf{Varianza:} $\Var{S_5} = 5/25 = 0.2$ min$^2$, $\sigma = \sqrt{0.2} \approx 0.447$ min.
  \\[3pt]
  \textbf{Probabilidad de esperar a lo m\'as $1$ min:}
  \[
    P(S_5 \leq 1) = P(N(1) \geq 5) = 1 - \sum_{k=0}^{4}\frac{5^k e^{-5}}{k!} \approx 0.560.
  \]
  Hay aproximadamente un $56\%$ de probabilidad de que las primeras cinco llamadas ocurran
  dentro del primer minuto.
\end{ejemplo}

\begin{center}
  \begin{tikzpicture}
    \begin{axis}[
        width=13cm, height=5.5cm,
        domain=0:10, samples=200,
        axis lines=left,
        xlabel={$t$}, ylabel={$f_{S_n}(t)$},
        xtick={0,1,2,3,4,5,6,7,8,9,10},
        ymin=0, ymax=1.05,
        legend style={at={(0.98,0.95)},anchor=north east,font=\small},
        tick label style={font=\small},
        label style={font=\small},
        grid=major, grid style={dotted, gray!40},
      ]
      \addplot[thick, color=sectionbg, domain=0:10] {exp(-x)};
      \addlegendentry{$n=1$}
      \addplot[thick, color=defborder, dashed, domain=0:10] {x*exp(-x)};
      \addlegendentry{$n=2$}
      \addplot[thick, color=noteborder, dash dot, domain=0:10] {x^2*exp(-x)/2};
      \addlegendentry{$n=3$}
      \addplot[thick, color=warnborder, densely dotted, domain=0:10]
      {x^4*exp(-x)/24};
      \addlegendentry{$n=5$}
    \end{axis}
  \end{tikzpicture}
\end{center}

\subsecheader{6.7 \textperiodcentered\ Procesos de Poisson bifurcados (\emph{thinning})}

\begin{proposicion}{Bifurcaci\'on}
  Sea $\{N(t)\}$ un proceso de Poisson de tasa $\lambda$. Cada evento se clasifica
  independientemente como tipo $1$ (prob.\ $p$) o tipo $2$ (prob.\ $1-p$). Entonces $N_1$ y
  $N_2$ son procesos de Poisson \textbf{independientes} con tasas $\lambda p$ y $\lambda(1-p)$.
\end{proposicion}

\begin{nota}
  \textbf{Sorpresa:} aunque comparten el mismo proceso original, los procesos resultantes son
  \emph{independientes}.
\end{nota}

\begin{ejemplo}{Tienda con compradores y curiosos}
  Llegan personas a una tienda como un proceso de Poisson de tasa $\lambda = 20$/hora.
  Cada persona compra con probabilidad $p = 0.3$, independientemente de las dem\'as.
  \\[3pt]
  \textbf{(1) Bifurcaci\'on:}
  \begin{itemize}
    \item Compradores: proceso de Poisson de tasa $\lambda p = 20 \cdot 0.3 = 6$/hora.
    \item No compradores: proceso de Poisson de tasa $\lambda(1-p) = 20 \cdot 0.7 = 14$/hora.
    \item Ambos procesos son \emph{independientes}.
  \end{itemize}
  \textbf{(2) N\'umero esperado de compradores en $2$ horas:} $6 \cdot 2 = 12$.
  \\
  \textbf{(3) Probabilidad de exactamente $5$ compradores en $1$ hora:}
  $P(N_1(1) = 5) = \frac{6^5 e^{-6}}{5!} \approx 0.161$.
  \\
  \textbf{(4) Independencia en acci\'on:} si en una hora hubo $14$ no-compradores (mucho m\'as
  de los $14$ esperados, $\Var = 14$), \emph{no} aporta informaci\'on sobre cu\'antos
  compradores hubo --- los procesos son completamente desacoplados.
\end{ejemplo}

\begin{cuidado}
  \textbf{Sorpresa de la independencia.} Aunque $N_1$ y $N_2$ se construyen del \emph{mismo}
  proceso original (claramente $N_1(t) + N_2(t) = N(t)$ en cada instante), resultan
  estad\'isticamente independientes. Esto contradice la intuici\'on de ``si suben los
  compradores, deben bajar los curiosos'' --- esa intuici\'on aplicar\'ia si fij\'aramos $N(t)$,
  pero no la fijamos. La elegancia del proceso de Poisson es que las fluctuaciones de cada
  tipo se compensan exactamente para que ambos sean Poisson independientes.
\end{cuidado}

\begin{center}
  \begin{tikzpicture}[scale=0.95, >=Latex]
    \draw[-Latex, thick] (-0.2,2) -- (8.5,2) node[right,font=\small] {$N(t)$};
    \foreach \x/\type in {1.0/1, 2.2/2, 3.0/1, 4.5/2, 5.7/1, 7.2/1} {
        \ifnum\type=1
          \fill[color=defborder] (\x,2) circle (2.5pt);
        \else
          \fill[color=noteborder] (\x,2) circle (2.5pt);
        \fi
      }
    \node[font=\small\bfseries, left] at (-0.3,2) {$\lambda$};
    \draw[-Latex, thick] (-0.2,0.5) -- (8.5,0.5) node[right,font=\small] {$N_1(t)$};
    \foreach \x in {1.0, 3.0, 5.7, 7.2} {
        \fill[color=defborder] (\x,0.5) circle (2.5pt);
      }
    \node[font=\small\bfseries, left, color=defborder] at (-0.3,0.5) {$\lambda p$};
    \draw[-Latex, thick] (-0.2,-1) -- (8.5,-1) node[right,font=\small] {$N_2(t)$};
    \foreach \x in {2.2, 4.5} {
        \fill[color=noteborder] (\x,-1) circle (2.5pt);
      }
    \node[font=\small\bfseries, left, color=noteborder] at (-0.3,-1) {$\lambda(1{-}p)$};
  \end{tikzpicture}
\end{center}

\subsecheader{6.8 \textperiodcentered\ Suma (superposici\'on) de procesos de Poisson}

\begin{proposicion}{Superposici\'on}
  Si $N_1, \ldots, N_k$ son procesos de Poisson independientes con tasas $\lambda_1, \ldots, \lambda_k$,
  entonces $N(t) = \sum_i N_i(t)$ es Poisson de tasa $\sum_i \lambda_i$.
\end{proposicion}

\begin{demostracion}
  Verificamos las tres propiedades del proceso de Poisson:
  \begin{enumerate}
    \item $N(0) = \sum N_i(0) = 0$.
    \item Cada $N_i$ tiene incrementos independientes, y entre procesos son independientes
          $\Rightarrow$ los incrementos de $N$ tambi\'en lo son.
    \item Para un intervalo de longitud $s$: $N_i(t+s)-N_i(t) \sim \mathrm{Poisson}(\lambda_i s)$.
          Por aditividad de Poissons independientes:
          $N(t+s) - N(t) \sim \mathrm{Poisson}(\sum \lambda_i\, s)$.
  \end{enumerate}
\end{demostracion}

\vspace{4pt}
\begin{center}
  \begin{tikzpicture}[scale=0.95, >=Latex]
    \draw[-Latex, thick] (-0.2,1.5) -- (8.5,1.5) node[right,font=\small] {$N_1(t)$};
    \foreach \x in {1.5, 3.5, 6.5} {
        \fill[color=defborder] (\x,1.5) circle (2.5pt);
      }
    \node[font=\small\bfseries, left, color=defborder] at (-0.3,1.5) {$\lambda_1$};

    \draw[-Latex, thick] (-0.2,0) -- (8.5,0) node[right,font=\small] {$N_2(t)$};
    \foreach \x in {0.8, 2.5, 4.2, 5.5, 7.5} {
        \fill[color=noteborder] (\x,0) circle (2.5pt);
      }
    \node[font=\small\bfseries, left, color=noteborder] at (-0.3,0) {$\lambda_2$};

    \draw[-Latex, thick] (-0.2,-1.5) -- (8.5,-1.5) node[right,font=\small] {$N(t)$};
    \foreach \x/\c in {0.8/n, 1.5/d, 2.5/n, 3.5/d, 4.2/n, 5.5/n, 6.5/d, 7.5/n} {
        \ifx\cn
          \fill[color=noteborder] (\x,-1.5) circle (2.5pt);
        \else
          \fill[color=defborder] (\x,-1.5) circle (2.5pt);
        \fi
      }
    \node[font=\small\bfseries, left] at (-0.3,-1.5) {$\lambda_1{+}\lambda_2$};

    \foreach \x in {1.5, 3.5, 6.5} {
        \draw[->, thin, color=defborder, opacity=0.5] (\x,1.35) -- (\x,-1.35);
      }
    \foreach \x in {0.8, 2.5, 4.2, 5.5, 7.5} {
        \draw[->, thin, color=noteborder, opacity=0.5] (\x,-0.15) -- (\x,-1.35);
      }
  \end{tikzpicture}
\end{center}
\begin{center}
  {\footnotesize\itshape Superposici\'on: dos procesos independientes se combinan en uno con
    tasa suma. Es la operaci\'on \emph{inversa} de la bifurcaci\'on.}
\end{center}

\begin{ejemplo}{Llamadas combinadas}
  Un call center recibe llamadas de dos sucursales independientes:
  sucursal $1$ a tasa $\lambda_1 = 3$/min, sucursal $2$ a tasa $\lambda_2 = 5$/min.
  \\[3pt]
  \textbf{(1) Tasa total:} el flujo combinado es Poisson de tasa $\lambda = 8$/min.\\
  \textbf{(2) Llamadas en $10$ min:} $N(10) \sim \mathrm{Poisson}(80)$, $\E{N(10)} = 80$.\\
  \textbf{(3) Tiempo entre llamadas (cualquier sucursal):} Exp$(8)$, media $1/8$ min $= 7.5$ s.\\
  \textbf{(4) Probabilidad de que la pr\'oxima llamada sea de la sucursal $1$:}
  $\lambda_1/(\lambda_1+\lambda_2) = 3/8 = 0.375$.
\end{ejemplo}

\subsecheader{6.9 \textperiodcentered\ Distribuci\'on condicional del primer evento}

\begin{proposicion}{$T_1 \mid N(t)=1$ es uniforme}
  En un proceso de Poisson de tasa $\lambda$, dado $N(t) = 1$:
  \[
    T_1 \mid N(t) = 1 \sim \mathrm{Unif}(0, t).
  \]
\end{proposicion}

\begin{demostracion}
  Para $0 \leq s \leq t$:
  \[
    P(T_1 \leq s \mid N(t) = 1) = \frac{P(T_1 \leq s,\; N(t) = 1)}{P(N(t) = 1)}.
  \]
  El evento $\{T_1 \leq s,\; N(t)=1\}$ es ``hubo exactamente $1$ evento en $(0,s]$ y $0$
  eventos en $(s,t]$''. Por incrementos independientes:
  \[
    P(T_1 \leq s, N(t) = 1) = P(N(s)=1)\cdot P(N(t)-N(s)=0)
    = \lambda s\, e^{-\lambda s} \cdot e^{-\lambda(t-s)} = \lambda s\, e^{-\lambda t}.
  \]
  Y $P(N(t)=1) = \lambda t\, e^{-\lambda t}$. Dividiendo:
  \[
    P(T_1 \leq s \mid N(t) = 1) = \frac{\lambda s\,e^{-\lambda t}}{\lambda t\,e^{-\lambda t}}
    = \frac{s}{t}.
  \]
  Esa es la FDA de $\mathrm{Unif}(0, t)$.
\end{demostracion}

\begin{ejemplo}{Ubicaci\'on del \'unico cliente}
  En una estaci\'on de servicio, los clientes llegan como Poisson($\lambda = 2$/hora).
  Sabemos que en $1$ hora lleg\'o exactamente $1$ cliente. \textquestiondown En qu\'e momento?
  \\[3pt]
  Por la proposici\'on, $T_1 \mid N(1) = 1 \sim \mathrm{Unif}(0, 1)$.
  \begin{itemize}
    \item Probabilidad de que llegara en los primeros $20$ minutos: $20/60 = 1/3 \approx 0.333$.
    \item Probabilidad de que llegara entre los minutos $30$ y $45$: $(45-30)/60 = 0.25$.
    \item Tiempo esperado de llegada: $\E{T_1 \mid N(1)=1} = 1/2$ hora $= 30$ min.
  \end{itemize}
  \textbf{Observaci\'on:} aunque \emph{a priori} $T_1 \sim \mathrm{Exp}(2)$ (con media $1/2 = 30$ min),
  al condicionar a que ocurri\'o exactamente uno en $(0, 1]$, el tiempo se vuelve
  uniformemente distribuido en el intervalo --- ya no exponencial.
\end{ejemplo}

\begin{nota}
  \textbf{Generalizaci\'on:} dado $N(t) = n$, los $n$ tiempos son los \emph{estad\'isticos de orden}
  de $n$ uniformes i.i.d.\ en $(0, t)$.
\end{nota}

\begin{center}
  \begin{tikzpicture}[scale=1, >=Latex]
    \draw[thick] (0,0) -- (8,0);
    \fill[color=sectionbg] (0,0) circle (2pt);
    \fill[color=sectionbg] (8,0) circle (2pt);
    \node[below, font=\small] at (0,-0.1) {$0$};
    \node[below, font=\small] at (8,-0.1) {$t$};
    \fill[fill=defbg, draw=defborder, opacity=0.7] (0,0.05) rectangle (8,0.6);
    \node[font=\small\itshape, color=defborder] at (4,0.32) {$f_{T_1\mid N(t)=1}(s) = 1/t$};
    \draw[decorate, decoration={brace, amplitude=4pt}, thick]
    (0,-0.4) -- (8,-0.4) node[midway, below=4pt, font=\small]
    {dado $N(t)=1$: $T_1 \sim \mathrm{Unif}(0,t)$};
  \end{tikzpicture}
\end{center}

\vspace{10pt}
\begin{tcolorbox}[
    colback=sectionbg!8, colframe=sectionbg,
    boxrule=0.8pt, arc=4pt,
    left=8pt, right=8pt, top=6pt, bottom=6pt,
    fonttitle=\bfseries\sffamily,
    title={Conexiones clave del m\'odulo},
  ]
  \small
  \begin{itemize}
    \item $\mathrm{Bin}(n,\lambda/n) \xrightarrow{n\to\infty} \mathrm{Poisson}(\lambda)$
          \hfill (eventos raros)
    \item $\mathrm{Poisson}(\lambda) \xrightarrow{\lambda\to\infty} \mathcal{N}(\lambda,\lambda)$
          \hfill (normalidad asint\'otica)
    \item Geom\'etrica $\Leftrightarrow$ Exponencial \hfill (p\'erdida de memoria)
    \item Tiempos entre llegadas Poisson $\sim$ Exp$(\lambda)$ i.i.d. \hfill (proceso de Poisson)
    \item $S_n \sim$ Erlang$(n, \lambda)$ \hfill (tiempo de la $n$-\'esima llegada)
    \item FDA $F_X$ determina la ley de $X$
  \end{itemize}
\end{tcolorbox}

\end{document}
